%%%%%%%%%%%%%%%%%%%%%%%%%%%%%%%%%%%%%%%%%%%%%%%%%%%%%%%%%%%%%%%%%%%%%%%%%%%%%%%
%  Laboratorio 3 - FFT con CMSIS-DSP en MCXN947
%  Estructura de carpetas (igual que en el TP1 y el TP2):
%     main.tex
%     informe-fcefyn.sty      -> copia de documentacion/TP2/
%     img/                    -> logos, capturas de Config Tools, de la GUI
%                                de MATLAB y del terminal serie
%
%  Imagenes que usa este documento (las que falten aparecen como un recuadro
%  rojo con el nombre esperado):
%     img/unc_logo.png, img/fcefyn_logo.jpg      (los mismos del TP2)
%     img/peripherals-lpuart-tp3.png             Config Tools > Peripherals >
%                                                LP_FLEXCOMM4 (LPUART 115200 8N1)
%     img/pins-tsi-tp3.png                       Config Tools > Pins: P1_3 como
%                                                TSI0_CH3
%     img/clocks-tsi-tp3.png                     Config Tools > Clocks: reloj del
%                                                TSI desde FRO_12M
%     img/terminal-trama.png                     Terminal serie con una trama
%                                                (cabecera, filas y END)
%     img/sim-fuga-1000.png, img/sim-fuga-1005.png
%                                                explorar_fft_tp3.m con f0 = 1000
%                                                y 1005 Hz (fs = 8 kS/s, N = 512)
%     img/gui-bypass-8k.png                      GUI en bypass, 1 kHz, 8 kS/s
%     img/gui-fft-8k-512.png                     GUI en FFT, 1 kHz, 8 kS/s, N=512
%     img/gui-fft-8k-1024.png                    idem N = 1024
%     img/gui-fft-8k-2048.png                    idem N = 2048
%     img/gui-fft-22k-512.png                    GUI en FFT, 1 kHz, 22 kS/s, N=512
%     img/gui-fft-48k-512.png                    GUI en FFT, 1 kHz, 48 kS/s, N=512
%     img/gui-alias-8k.png                       5 kHz a 8 kS/s (aliasing)
%     img/gui-alias-16k.png                      5 kHz a 16 kS/s
%     img/gui-cuadrada-8k.png                    cuadrada de 500 Hz a 8 kS/s
%
%  El Anexo incluye ../../TP3/source/TP3.c y ../../TP3/matlab/recibir_tp3.m,
%  con la ruta relativa a esta carpeta (si se compila fuera del repositorio
%  aparece una marca roja en su lugar).
%
%  Todo lo que aparece entre corchetes rojos (\completar{...}) es contenido
%  que falta llenar. Para la entrega, descomentar \entregafinal: el documento
%  avisara en el log si quedo alguna marca sin resolver.
%%%%%%%%%%%%%%%%%%%%%%%%%%%%%%%%%%%%%%%%%%%%%%%%%%%%%%%%%%%%%%%%%%%%%%%%%%%%%%%
\documentclass[11pt,a4paper]{article}
\usepackage{informe-fcefyn}

% \entregafinal

% ------------------------- Datos de la caratula ------------------------------
\materia{Procesamiento Digital de Señales}
\carrera{Ingeniería en Computación}
\subtitulo{Laboratorio 3}
\tema{FFT con CMSIS-DSP en MCXN947}
\plataforma{FRDM-MCXN947 \textperiodcentered{} MCUXpresso IDE \textperiodcentered{} CMSIS-DSP \textperiodcentered{} MATLAB}
\autor{Costamagna, Matias Javier}{44.512.437}{matias.costamagna@mi.unc.edu.ar}
\autor{Moreyra, Julián}{45.214.537}{julian.moreyra@mi.unc.edu.ar}
\autor{Rivera, Mariano}{45.432.921}{luismarianorivera.25@mi.unc.edu.ar}
\autor{Rojas, Facundo Nicolás}{44.629.222}{facundo.rojas.222@mi.unc.edu.ar}
\docentes{Rossi, Roberto Raúl \quad\textperiodcentered\quad Parlanti, Gustavo Fernando \quad\textperiodcentered\quad Molina, Germán Alcides \quad\textperiodcentered\quad Ferreyra, Pablo Alejandro}
\fechaentrega{\today}
\logoizquierdo{img/unc_logo.png}
\logoderecho{img/fcefyn_logo.jpg}

% ---------------- Parametros del diseno (usados en las tablas) ---------------
% Igual que en el TP1 y el TP2: CTIMER0 a 150 MHz, match 3 en toggle.
\frecuenciatimer{150000000}
\divisordisparo{2}

% Velocidad de la UART [bit/s] y bits por caracter (8N1: inicio + 8 + parada)
\newcommand{\baudios}{115200}
\newcommand{\bitsporbyte}{10}

% Tipos y funciones propios de este trabajo en los listados de codigo
\lstset{morekeywords=[2]{arm_cfft_instance_f32,tsi_selfCap_config_t,size_t}}

% \filaresol{fs nominal [Hz]}{MR}: fs real, resolucion fs/N y duracion N/fs
% de la ventana para N = 512, 1024 y 2048.
\newcommand{\filaresol}[2]{%
  \xdef\fsrealtmp{\fpeval{150000000/(2*(#2+1))}}%
  \xdef\dfa{\fpeval{round(\fsrealtmp/512,2)}}%
  \xdef\dfb{\fpeval{round(\fsrealtmp/1024,2)}}%
  \xdef\dfc{\fpeval{round(\fsrealtmp/2048,2)}}%
  \xdef\tva{\fpeval{round(1000*512/\fsrealtmp,1)}}%
  \xdef\tvb{\fpeval{round(1000*1024/\fsrealtmp,1)}}%
  \xdef\tvc{\fpeval{round(1000*2048/\fsrealtmp,1)}}%
  \qty{#1}{\hertz} & \num{\dfa} & \num{\dfb} & \num{\dfc} &
  \num{\tva} & \num{\tvb} & \num{\tvc}\\}

% \filacomp{N}: productos complejos de la DFT directa y de la FFT radix-2.
\newcommand{\filacomp}[1]{%
  \xdef\cdirtmp{\fpeval{(#1)*(#1)}}%
  \xdef\cffttmp{\fpeval{round((#1)/2*ln(#1)/ln(2),0)}}%
  \xdef\crattmp{\fpeval{round((#1)*(#1)/((#1)/2*ln(#1)/ln(2)),1)}}%
  \xdef\etapastmp{\fpeval{round(ln(#1)/ln(2),0)}}%
  \num{#1} & \num{\etapastmp} & \num{\cdirtmp} & \num{\cffttmp} &
  \num{\crattmp}\\}

% \filatx{N}{bytes de la trama}: tiempo de transmision, de adquisicion y
% periodo de refresco del espectro a 8 kS/s y a 48 kS/s.
\newcommand{\filatx}[2]{%
  \xdef\txtmp{\fpeval{round((#2)*\bitsporbyte/\baudios,2)}}%
  \xdef\acqatmp{\fpeval{round(1000*(#1)/8000,1)}}%
  \xdef\acqbtmp{\fpeval{round(1000*(#1)/47984.64,1)}}%
  \xdef\peratmp{\fpeval{round((#1)/8000+(#2)*\bitsporbyte/\baudios,2)}}%
  \xdef\perbtmp{\fpeval{round((#1)/47984.64+(#2)*\bitsporbyte/\baudios,2)}}%
  \num{#1} & \num{#2} & \num{\txtmp} & \num{\acqatmp} & \num{\acqbtmp} &
  \num{\peratmp} & \num{\perbtmp}\\}

\begin{document}

\caratula
\setcounter{tocdepth}{2}
\tableofcontents
\clearpage

%%%%%%%%%%%%%%%%%%%%%%%%%%%%%%%%%%%%%%%%%%%%%%%%%%%%%%%%%%%%%%%%%%%%%%%%%%%%%%%
\section{Introducción}
%%%%%%%%%%%%%%%%%%%%%%%%%%%%%%%%%%%%%%%%%%%%%%%%%%%%%%%%%%%%%%%%%%%%%%%%%%%%%%%

\begin{consigna}[Laboratorio 3 \textemdash{} Transformada Rápida de Fourier]
\textbf{Objetivo.} Realizar la implementación de la Transformada Rápida de
Fourier (FFT) de 512 muestras.

\textbf{Descripción.} Realizar un programa aplicativo que sea capaz de aplicar
la FFT a un buffer de memoria de diferente tamaño (por ejemplo, 512, 1024 y
2048 muestras), que es adquirido con el Laboratorio~1 teniendo en cuenta
además los cambios de velocidades de muestreo. Con una de las teclas de la
placa de evaluación se habilitará la aplicación de la FFT o se hará
\emph{bypass} de la misma a un buffer de salida distinto del buffer de
entrada. Transmitir el buffer (resultado de la FFT) por el puerto serie (SCI)
y graficar en la PC con alguna aplicación (por ejemplo, una GUI de MATLAB) el
espectro obtenido. Sacar conclusiones respecto de la implementación de la FFT
cuando cambia la frecuencia de muestreo o el tamaño de la ventana de
muestras.
\end{consigna}

\subsection{Requerimientos}

De la consigna se desprenden los siguientes requerimientos, que se usan como
guía para el diseño y como lista de verificación en la
sección~\ref{sec:verificacion}:

\begin{enumerate}[label=\textbf{R\arabic*},leftmargin=2.2em]
  \item \label{req:fs} Adquirir la señal con la aplicación del
        Laboratorio~1, con sus cinco velocidades de muestreo
        (\qtylist{8;16;22;44;48}{\kilo\sps}) seleccionables durante la
        ejecución.
  \item \label{req:n} Guardar las muestras en un buffer de entrada de
        tamaño variable: 512, 1024 o 2048 muestras.
  \item \label{req:fft} Aplicar la FFT al buffer de entrada.
  \item \label{req:tecla} Con una tecla, habilitar la FFT o hacer
        \emph{bypass} hacia un buffer de salida distinto del de entrada.
  \item \label{req:sci} Transmitir el buffer de salida por el puerto serie.
  \item \label{req:gui} Graficar en la PC, con una aplicación de MATLAB, el
        espectro recibido.
  \item \label{req:concl} Sacar conclusiones sobre el efecto de la frecuencia
        de muestreo y del tamaño de la ventana en la FFT.
\end{enumerate}

\begin{nota}[Decisiones de implementación respecto de la consigna]
\begin{itemize}[nosep]
  \item El tamaño de la ventana se elige durante la ejecución, sin
        recompilar: el touchpad capacitivo E1 recorre 512, 1024 y 2048
        muestras. SW3 conserva la selección circular de la frecuencia de
        muestreo de los laboratorios anteriores y SW2 alterna entre FFT y
        \emph{bypass}, como pide la consigna.
  \item El procesamiento es \emph{por bloques}: la rutina del ADC llena el
        buffer de entrada y el lazo principal calcula la FFT y transmite el
        resultado. Mientras tanto la adquisición sigue disparándose por
        hardware, pero las muestras se descartan hasta que el bloque termina
        de enviarse. Cada espectro corresponde entonces a una ventana de
        $N$ muestras consecutivas, y entre dos espectros hay un hueco. La
        sección~\ref{sec:tiempos} muestra que con el puerto serie a
        \qty{115200}{\bit\per\second} no es posible transmitir espectros
        en forma continua.
  \item La FFT se calcula en punto flotante de simple precisión con
        \texttt{arm\_cfft\_f32()} de CMSIS-DSP, aprovechando la FPU del
        Cortex-M33 (sección~\ref{sec:formato}).
  \item En \emph{bypass} el buffer de salida contiene las muestras de la
        ventana, sin transformar, y la GUI las grafica en el dominio del
        tiempo. Así el mismo enlace sirve para ver la señal y su espectro.
  \item El TP3 no usa el DAC: la salida es el puerto serie. El LED RGB
        conserva el color de cada frecuencia de muestreo; el modo, el tamaño
        de la ventana y la $f_s$ de cada bloque viajan en la cabecera de la
        trama y la GUI los muestra en el título del gráfico.
\end{itemize}
\end{nota}

%%%%%%%%%%%%%%%%%%%%%%%%%%%%%%%%%%%%%%%%%%%%%%%%%%%%%%%%%%%%%%%%%%%%%%%%%%%%%%%
\ejercicio{Configuración del ambiente}\label{sec:ambiente}
%%%%%%%%%%%%%%%%%%%%%%%%%%%%%%%%%%%%%%%%%%%%%%%%%%%%%%%%%%%%%%%%%%%%%%%%%%%%%%%

\subsection{Herramientas utilizadas}

\begin{table}[H]
\centering
\caption{Hardware y software del entorno de trabajo.}
\label{tab:herramientas}
\begin{tabular}{@{}l l@{}}
\toprule
\textbf{Elemento} & \textbf{Versión / modelo}\\
\midrule
Placa de evaluación     & FRDM-MCXN947 (Cortex-M33 con FPU, núcleo \texttt{cm33\_core0})\\
IDE                     & MCUXpresso IDE v25.6.136\\
SDK                     & \texttt{SDK\_26.06.00\_FRDM-MCXN947}\\
Herramientas de config. & MCUXpresso Config Tools (Clocks, Pins, Peripherals)\\
Proyecto base           & Proyecto del Laboratorio 2 (\archivo{TP2})\\
Drivers                 & SDK: \texttt{fsl\_lpuart}, \texttt{fsl\_tsi\_v6}, \texttt{fsl\_lpadc}, \texttt{fsl\_ctimer}\\
Biblioteca DSP          & CMSIS-DSP (\texttt{arm\_cfft\_f32})\\
Visualización           & MATLAB (función \texttt{serialport}, R2019b o posterior)\\
Enlace serie            & Puerto COM virtual del MCU-Link integrado (USB)\\
Generador de señales    & \\
Osciloscopio            & \\
\bottomrule
\end{tabular}
\end{table}

\subsection{Creación del proyecto}

El proyecto \archivo{TP3} parte del proyecto del Laboratorio 2: conserva la
cadena CTIMER0 $\rightarrow$ INPUTMUX $\rightarrow$ LPADC0, la selección
circular de la frecuencia de muestreo con SW3, el color del LED según la
$f_s$ y la lectura del touchpad sin interrupciones. Sobre esa base:

\begin{enumerate}[nosep]
  \item Se quitó el DAC0 y el filtrado muestra a muestra: la rutina del ADC
        sólo guarda muestras, y el procesamiento pasa al lazo principal.
  \item Se mantuvo el código fuente de CMSIS-DSP (carpeta
        \archivo{CMSIS/DSP}), que además de los filtros incluye las
        funciones de transformadas (\archivo{TransformFunctions}) y las
        tablas precalculadas (\archivo{CommonTables}):
        \texttt{arm\_cfft\_f32()}, \texttt{arm\_bitreversal\_32()} y las
        instancias constantes \texttt{arm\_cfft\_sR\_f32\_len512},
        \texttt{\_len1024} y \texttt{\_len2048} de
        \archivo{arm_const_structs.h}.
  \item Se configuró en Config Tools la instancia \texttt{LP\_FLEXCOMM4}
        como LPUART (sección~\ref{sec:cambios-perif}), que es la que la placa
        conecta al puerto COM virtual del MCU-Link.
  \item Se agregó el componente \texttt{tsi\_v6} del SDK (driver
        \archivo{fsl_tsi_v6}) para leer el touchpad, igual que en el
        Laboratorio 2.
  \item La lógica de la aplicación está en \archivo{source/TP3.c} y la GUI y
        las herramientas de prueba de MATLAB en \archivo{TP3/matlab/}
        (sección~\ref{sec:matlab}).
\end{enumerate}

\subsection{Relojes y pines}

Los relojes son los del Laboratorio~2: PLL0 a \qty{150}{\mega\hertz} para el
núcleo y el CTIMER0, y FRO\_HF de \qty{48}{\mega\hertz} para el ADC0. Se
agrega el FRO de \qty{12}{\mega\hertz} como reloj funcional del
LP\_FLEXCOMM4 (la UART) y del TSI. La tabla~\ref{tab:pines} resume la
asignación de pines.

\begin{table}[H]
\centering
\caption{Asignación de pines. Respecto del Laboratorio 2 desaparece la salida
del DAC y se agrega la UART.}
\label{tab:pines}
\footnotesize
\begin{tabular}{@{}l l l l@{}}
\toprule
\textbf{Función} & \textbf{Pin} & \textbf{Periférico} & \textbf{Conector / observación}\\
\midrule
Entrada analógica          & \pin{ADC0\_A0}          & LPADC0 canal A0    & \texttt{J4[2]} (Arduino A0)\\
Disparo del ADC            & \pin{P0\_19}            & CT0\_MAT3          & \texttt{J2[13]}, punto de prueba\\
LED rojo / verde / azul    & \pin{P0\_10}, \pin{P0\_27}, \pin{P1\_2} & GPIO0, GPIO1 & Activos en bajo\\
Tecla SW3 (velocidad)      & \pin{P0\_6}             & GPIO0, canal IRQ 0 & Flanco descendente, pull-up interno\\
Tecla SW2 (FFT / \emph{bypass}) & \pin{P0\_23}       & GPIO0, canal IRQ 1 & Flanco descendente\\
Touchpad E1 (tamaño $N$)   & \pin{P1\_3}             & TSI0, canal 3      & Self-cap, disparo por software\\
UART RXD / TXD             & \pin{P1\_8}, \pin{P1\_9} & LP\_FLEXCOMM4      & COM virtual del MCU-Link\\
\bottomrule
\end{tabular}
\end{table}

% ATENCION: en el commit 8b11c19 el proyecto TP3 no tiene asignado P1_3 como
% TSI0_CH3 en pin_mux.c ni el reloj del TSI (TSICLKSEL = FRO_12M, TSICLKDIV)
% en clock_config.c. Sin eso init_touch() no completa el barrido de
% calibracion, deja touch_ok en false y el tamano queda fijo en 512. Hay que
% replicar en Config Tools lo hecho en el TP2 (ver TP2/board/pin_mux.c y
% TP2/board/clock_config.c) y borrar el \completar de abajo.
\completar{Verificar que el TP3 tenga en Config Tools \pin{P1\_3} como
\texttt{TSI0\_CH3} y el reloj del TSI desde FRO\_12M, como en el TP2: en el
estado actual del repositorio faltan, y sin ellos el touchpad queda
deshabilitado y $N$ fijo en 512.}

\captura{img/pins-tsi-tp3.png}{Config Tools \textemdash{} Pins}{\pin{P1\_3}
asignado a \texttt{TSI0\_CH3} (touchpad E1), con el buffer de entrada
digital deshabilitado.}{fig:pins-tsi}

\captura{img/clocks-tsi-tp3.png}{Config Tools \textemdash{} Clocks}{Reloj
del TSI tomado del FRO de \qty{12}{\mega\hertz}.}{fig:clocks-tsi}

\subsection{Cambios en los periféricos}\label{sec:cambios-perif}

La tabla~\ref{tab:cambios} resume lo que cambió respecto del Laboratorio 2.
El resto de la configuración (CTIMER0 en modo toggle sobre el match 3,
disparo del ADC por INPUTMUX, LPADC0 en \qty{16}{\bit} con referencia
\texttt{kLPADC\_ReferenceVoltageAlt3}, interrupciones de las teclas en las
líneas 0 y 1 del GPIO0) se mantiene.

\begin{table}[H]
\centering
\caption{Diferencias de configuración respecto del Laboratorio 2.}
\label{tab:cambios}
\footnotesize
\begin{tabular}{@{}l p{5.0cm} p{6.4cm}@{}}
\toprule
\textbf{Elemento} & \textbf{Laboratorio 2} & \textbf{Laboratorio 3}\\
\midrule
DAC0                  & \texttt{init\_dac()}, una muestra por interrupción
                      & No se usa\\
VREF                  & Encendido por \texttt{init\_dac()}
                        (\texttt{SPC0->ACTIVE\_CFG1})
                      & Encendido explícito con
                        \texttt{SPC\_EnableActive}\allowbreak\texttt{ModeAnalogModules()}
                        antes de los periféricos\\
LP\_FLEXCOMM4         & Consola de depuración (\texttt{PRINTF})
                      & LPUART de Config Tools, \num{115200}~bit/s, 8N1, sólo
                        transmisión, por \emph{polling}\\
Tecla SW2             & Selección de filtro & FFT / \emph{bypass}\\
Touchpad E1           & Habilita el filtro  & Elige $N$: 512 $\rightarrow$ 1024
                                              $\rightarrow$ 2048 $\rightarrow$ 512\\
LED RGB               & Color según la $f_s$; parpadeo según el filtro
                      & Color según la $f_s$\\
Procesamiento         & Muestra a muestra, en la rutina del ADC
                      & Por bloques de $N$ muestras, en el lazo principal\\
Buffers               & \texttt{buffer\_in[512]} y \texttt{buffer\_out[512]}
                        (q15, circulares)
                      & \texttt{buffer\_in[2048]} (q15) y
                        \texttt{buffer\_out[4096]} (float32, complejo
                        intercalado), lineales\\
CMSIS-DSP             & \texttt{arm\_fir\_q15()}
                      & \texttt{arm\_cfft\_f32()} y tablas de
                        \archivo{arm_const_structs.h}\\
Lazo principal        & Touchpad y LED
                      & Touchpad, FFT y transmisión\\
\bottomrule
\end{tabular}
\end{table}

\captura{img/peripherals-lpuart-tp3.png}{Config Tools \textemdash{}
LP\_FLEXCOMM4}{Instancia LPUART en modo \emph{polling}: \num{115200}~bit/s,
8 bits de datos, sin paridad, 1 bit de parada, sólo transmisión.}{fig:lpuart}

\subsubsection{Orden de inicialización y referencia del ADC}

El LPADC0 usa la referencia \texttt{kLPADC\_ReferenceVoltageAlt3}, que
proviene del módulo VREF. En el Laboratorio 2 el VREF quedaba encendido como
efecto lateral de \texttt{init\_dac()}, que escribía
\texttt{SPC0->ACTIVE\_CFG1 |= 0x11} para alimentar el DAC0 y el VREF. Al
quitar el DAC en el TP3, el VREF quedó apagado y el ADC dejó de entregar
bloques válidos: no salía ninguna trama por la UART. Por eso
\texttt{main()} enciende el VREF en forma explícita con
\begin{center}
\texttt{SPC\_EnableActiveModeAnalogModules(SPC0, kSPC\_controlVref);}
\end{center}
\emph{antes} de llamar a \texttt{BOARD\_InitBootPeripherals()}, que calibra el
ADC y arranca el CTIMER0. Así la calibración y la primera conversión ya se hacen
con la referencia estable.

\subsubsection{Prioridades de interrupción}

Igual que en los laboratorios anteriores, \texttt{ADC0\_IRQn},
\texttt{GPIO00\_IRQn} y \texttt{GPIO01\_IRQn} quedan con la prioridad por
defecto (0), así que ninguna desaloja a otra y las variables que comparten
no pueden cambiar en medio de una rutina. El SysTick, que sólo cuenta
milisegundos para el touchpad, queda con la prioridad mínima
(\texttt{SysTick\_Config()} la fija así).

A diferencia del Laboratorio 2, el trabajo pesado (la FFT y la transmisión)
no está en ninguna rutina de interrupción sino en el lazo principal, que
corre en modo \emph{thread} y puede ser interrumpido en cualquier momento.
La rutina del ADC, que es corta, sigue atendiendo cada conversión durante
los segundos que dura una transmisión, y por eso la FIFO del ADC nunca se
desborda aunque las muestras no se guarden.

%%%%%%%%%%%%%%%%%%%%%%%%%%%%%%%%%%%%%%%%%%%%%%%%%%%%%%%%%%%%%%%%%%%%%%%%%%%%%%%
\ejercicio{La Transformada Rápida de Fourier}\label{sec:teoria}
%%%%%%%%%%%%%%%%%%%%%%%%%%%%%%%%%%%%%%%%%%%%%%%%%%%%%%%%%%%%%%%%%%%%%%%%%%%%%%%

\subsection{Transformada discreta de Fourier}

La transformada discreta de Fourier (DFT) de $N$ puntos de una secuencia
$x[n]$, $n = 0, \dots, N-1$, es
\begin{equation}
  X[k] = \sum_{n=0}^{N-1} x[n]\, W_N^{kn},
  \qquad W_N = e^{-j 2\pi / N},
  \qquad k = 0, \dots, N-1.
  \label{eq:dft}
\end{equation}
Si $x[n]$ son muestras tomadas a $f_s$, el coeficiente $X[k]$ corresponde a
la frecuencia $f_k = k\, f_s / N$, así que la DFT evalúa el espectro en una
grilla de frecuencias equiespaciadas con separación
\begin{equation}
  \Delta f = \frac{f_s}{N} = \frac{1}{T_v},
  \qquad T_v = \frac{N}{f_s},
  \label{eq:resolucion}
\end{equation}
donde $T_v$ es la duración de la ventana analizada. La
ecuación~\eqref{eq:resolucion} es la que gobierna todo el trabajo: la
resolución en frecuencia es la inversa de la duración de la ventana, y no
depende de $f_s$ y de $N$ por separado sino de su cociente.

Cuando la entrada es real, como la del ADC, los coeficientes cumplen
$X[N-k] = X^*[k]$: la mitad superior del espectro es la imagen conjugada de
la inferior y la información está en los $N/2+1$ coeficientes que van de
\qty{0}{\hertz} a $f_s/2$. Los $k > N/2$ corresponden a frecuencias
negativas.

\subsection{Algoritmo de la FFT}

Calcular la ecuación~\eqref{eq:dft} en forma directa requiere $N$ productos
complejos por coeficiente, es decir, $N^2$ en total. La FFT obtiene el mismo
resultado aprovechando la periodicidad y la simetría de $W_N^{kn}$. En la
versión de diezmado en el tiempo (radix-2) la secuencia se separa en sus
muestras pares e impares, $g[m] = x[2m]$ y $h[m] = x[2m+1]$, cuyas DFT de
$N/2$ puntos son $G[k]$ y $H[k]$. Entonces
\begin{equation}
  X[k] = G[k] + W_N^{k} H[k], \qquad
  X[k + N/2] = G[k] - W_N^{k} H[k], \qquad
  k = 0, \dots, N/2 - 1,
  \label{eq:mariposa}
\end{equation}
porque $W_N^{k+N/2} = -W_N^{k}$. Cada par de la
ecuación~\eqref{eq:mariposa} es una \emph{mariposa}: un producto complejo
por el factor de giro (\emph{twiddle}) $W_N^k$ y dos sumas complejas.
Aplicando la división recursivamente hasta DFT de 2 puntos quedan
$\log_2 N$ etapas de $N/2$ mariposas cada una, con un total de
$(N/2)\log_2 N$ productos complejos. El precio es el orden: si el cálculo se
hace en el lugar (\emph{in place}), una de las dos secuencias, la entrada o
la salida, queda en orden de \emph{bits invertidos} (el índice $n$ escrito
en binario y leído al revés) y hay que reordenarla.

\begin{table}[H]
\centering
\caption{Productos complejos de la DFT directa y de la FFT radix-2 para los
tres tamaños de ventana.}
\label{tab:complejidad}
\begin{tabular}{@{}r r r r r@{}}
\toprule
\textbf{$N$} & \textbf{Etapas} & \textbf{DFT: $N^2$} &
\textbf{FFT: $(N/2)\log_2 N$} & \textbf{Relación}\\
\midrule
\filacomp{512}
\filacomp{1024}
\filacomp{2048}
\bottomrule
\end{tabular}
\end{table}

La ventaja crece con $N$: al duplicar la ventana la DFT directa cuesta
cuatro veces más y la FFT apenas algo más del doble. Para $N = 2048$ la DFT
directa necesitaría más de cuatro millones de productos complejos, que a
\qty{150}{\mega\hertz} y con varios ciclos por producto serían más de
\qty{100}{\milli\second}; la FFT lo resuelve con unos \num{11000}.

\subsection{Implementación en CMSIS-DSP}\label{sec:cmsis-fft}

\texttt{arm\_cfft\_f32()} calcula la FFT compleja de $N$ puntos en punto
flotante de simple precisión. Recibe un arreglo de $2N$ valores
\texttt{float32\_t} con la parte real y la imaginaria intercaladas,
$\{\mathrm{Re}\,x[0], \mathrm{Im}\,x[0], \mathrm{Re}\,x[1], \dots\}$, y
devuelve el resultado en el mismo arreglo. Sus parámetros son:

\begin{itemize}[nosep]
  \item \texttt{S}: puntero a una instancia \texttt{arm\_cfft\_instance\_f32}
        con la longitud $N$, la tabla de factores de giro y la tabla de
        reordenamiento. CMSIS-DSP las trae precalculadas en flash como
        constantes (\texttt{arm\_cfft\_sR\_f32\_len512}, \dots), así que no
        hay que inicializar nada en tiempo de ejecución.
  \item \texttt{ifftFlag}: 0 para la transformada directa, 1 para la
        inversa.
  \item \texttt{bitReverseFlag}: 1 para reordenar la salida y obtener los
        coeficientes en orden natural ($X[0], X[1], \dots$).
\end{itemize}

Internamente no usa mariposas radix-2 sino radix-8, que combinan ocho
entradas por mariposa y reducen la cantidad de etapas a $\log_8 N$ y la de
multiplicaciones por factores de giro. Como $512 = 8^3$, la FFT de 512 puntos
se resuelve sólo con radix-8 (\texttt{arm\_radix8\_butterfly\_f32()}). Para
1024 y 2048 agrega primero una etapa radix-2 o radix-4 que divide el problema
en dos o cuatro FFT de 512 (\texttt{arm\_cfft\_radix8by2\_f32()} y
\texttt{arm\_cfft\_radix8by4\_f32()}). Al final, si
\texttt{bitReverseFlag} vale 1, \texttt{arm\_bitreversal\_32()} intercambia
los pares de posiciones indicados en la tabla de reordenamiento. La
tabla~\ref{tab:tablas} muestra la memoria que ocupan esas tablas.

\begin{table}[H]
\centering
\caption{Tablas precalculadas de CMSIS-DSP para cada tamaño (en flash).}
\label{tab:tablas}
\small
\begin{tabular}{@{}r l r l r@{}}
\toprule
\textbf{$N$} & \textbf{Factores de giro} & \textbf{Bytes} &
\textbf{Reordenamiento} & \textbf{Bytes}\\
\midrule
512  & \texttt{twiddleCoef\_512}, \num{1024} float  & \num{4096}  &
       \texttt{armBitRevIndexTable512}, 448 \texttt{uint16}   & 896\\
1024 & \texttt{twiddleCoef\_1024}, \num{2048} float & \num{8192}  &
       \texttt{armBitRevIndexTable1024}, \num{1800} \texttt{uint16} & \num{3600}\\
2048 & \texttt{twiddleCoef\_2048}, \num{4096} float & \num{16384} &
       \texttt{armBitRevIndexTable2048}, \num{3808} \texttt{uint16} & \num{7616}\\
\bottomrule
\end{tabular}
\end{table}

CMSIS-DSP también ofrece \texttt{arm\_rfft\_fast\_f32()}, que aprovecha que
la entrada es real: calcula una FFT compleja de $N/2$ puntos y la separa en
los $N/2+1$ coeficientes no redundantes, con aproximadamente la mitad de
operaciones y de memoria. Se eligió la transformada compleja porque su
salida es directamente la DFT completa de la ecuación~\eqref{eq:dft}, sin
el formato empaquetado de la real, y porque con una entrada real basta con
poner en cero la parte imaginaria. La sección~\ref{sec:tiempos} muestra que
el costo de cálculo es despreciable frente al de la transmisión, así que el
ahorro no cambia el comportamiento del sistema.

\subsection{Interpretación de la amplitud}\label{sec:amplitud}

\texttt{arm\_cfft\_f32()} no normaliza la transformada directa. Para una
senoidal de amplitud $A$ (en las unidades normalizadas de la
sección~\ref{sec:formato}) cuya frecuencia coincide exactamente con un
coeficiente $k_0$, la ecuación~\eqref{eq:dft} da $|X[k_0]| = |X[N-k_0]| =
A\,N/2$: la amplitud se reparte entre la frecuencia positiva y su imagen
negativa, y cada una crece con $N$. Para que el gráfico no dependa del
tamaño de la ventana, la GUI dibuja el espectro de un solo lado,
\begin{equation}
  A[0] = \frac{|X[0]|}{N}, \qquad
  A[k] = \frac{2\,|X[k]|}{N}, \quad 1 \le k \le \frac{N}{2} - 1, \qquad
  A[N/2] = \frac{|X[N/2]|}{N},
  \label{eq:amplitud}
\end{equation}
con lo que una senoidal de amplitud $A$ aparece con altura $A$ y la
componente continua con su valor medio. Se calcula en la PC y no en la
placa para enviar el resultado exacto de \texttt{arm\_cfft\_f32()}.

\subsection{Fuga espectral y ventana}\label{sec:fuga}

La DFT trata la ventana de $N$ muestras como un período de una señal
periódica. Si la señal tiene un número entero de ciclos en la ventana, su
frecuencia cae exactamente sobre un coeficiente y el espectro tiene un único
valor no nulo. Si no, la extensión periódica tiene un salto en los bordes y
la energía se reparte entre todos los coeficientes: es la \emph{fuga
espectral} (\emph{leakage}). Equivale a multiplicar la señal por una
ventana rectangular, cuya transformada es un núcleo de Dirichlet con lóbulo
principal de ancho $2\,\Delta f$ y lóbulos laterales que empiezan en
\qty{-13.3}{\deci\bel} y decaen lentamente.

La fuga tiene dos consecuencias prácticas:
\begin{itemize}[nosep]
  \item \textbf{Error de amplitud} (\emph{scalloping}): si la frecuencia cae
        entre dos coeficientes, el pico es menor que $A$. En el peor caso,
        a medio coeficiente, la pérdida es de \qty{3.92}{\deci\bel}
        ($A$ se ve como $0{,}64\,A$).
  \item \textbf{Enmascaramiento}: los lóbulos laterales de un tono fuerte
        pueden tapar a un tono débil cercano.
\end{itemize}

El script \archivo{explorar_fft_tp3.m} lo ilustra con $f_s =
\qty{8}{\kilo\sps}$ y $N = 512$ ($\Delta f = \qty{15.625}{\hertz}$). Una
senoidal de \qty{1000}{\hertz} tiene exactamente 64 ciclos en la ventana:
cae en el coeficiente 64 y la amplitud se lee sin error. Con
\qty{1005}{\hertz} hay $64{,}32$ ciclos: el pico queda en
\qty{1000}{\hertz} con amplitud $0{,}84\,A$, el coeficiente vecino
(\qty{1015.6}{\hertz}) recibe $0{,}40\,A$ y los demás no se anulan
(tabla~\ref{tab:fuga}).

\begin{table}[H]
\centering
\caption{Amplitud de un solo lado, ecuación~\eqref{eq:amplitud}, alrededor
del pico para una senoidal de $A = 0{,}5$ cuantizada a q15, con $f_s =
\qty{8}{\kilo\sps}$ y $N = 512$ (calculada con la DFT de la
ecuación~\eqref{eq:dft}).}
\label{tab:fuga}
\begin{tabular}{@{}r r r r@{}}
\toprule
\textbf{$k$} & \textbf{$f_k$ [Hz]} & \textbf{$A[k]$ con $f_0 = \qty{1000}{\hertz}$} &
\textbf{$A[k]$ con $f_0 = \qty{1005}{\hertz}$}\\
\midrule
62 & \num{968.75}   & \num{0.0000} & \num{0.0575}\\
63 & \num{984.375}  & \num{0.0000} & \num{0.1014}\\
64 & \num{1000}     & \num{0.5000} & \num{0.4195}\\
65 & \num{1015.625} & \num{0.0000} & \num{0.1980}\\
66 & \num{1031.25}  & \num{0.0000} & \num{0.0804}\\
\bottomrule
\end{tabular}
\end{table}

La implementación no aplica ninguna ventana además de la rectangular, para
que el espectro enviado sea exactamente la FFT del buffer. Una ventana de
Hann (\texttt{arm\_hanning\_f32()} en CMSIS-DSP) bajaría los lóbulos
laterales a \qty{-31.5}{\deci\bel} y el error de amplitud máximo a
\qty{1.42}{\deci\bel}, a cambio de duplicar el ancho del lóbulo principal
(sección~\ref{sec:mejoras}).

%%%%%%%%%%%%%%%%%%%%%%%%%%%%%%%%%%%%%%%%%%%%%%%%%%%%%%%%%%%%%%%%%%%%%%%%%%%%%%%
\ejercicio{Diseño de la aplicación}\label{sec:diseno}
%%%%%%%%%%%%%%%%%%%%%%%%%%%%%%%%%%%%%%%%%%%%%%%%%%%%%%%%%%%%%%%%%%%%%%%%%%%%%%%

\subsection{Arquitectura general}

La figura~\ref{fig:bloques} muestra las partes de la aplicación. La
adquisición es la del Laboratorio~1: el CTIMER0 dispara cada conversión a
través del INPUTMUX y el LPADC0 deja el resultado en su FIFO. La rutina de
interrupción del ADC lo convierte a q15 y lo guarda en \texttt{buffer\_in}
hasta completar un bloque de $N$ muestras; entonces levanta la bandera
\texttt{buffer\_listo}. El lazo principal ve la bandera, convierte el bloque
a punto flotante en \texttt{buffer\_out}, le aplica la FFT (o no, en
\emph{bypass}) y lo transmite por la UART al puerto COM virtual del
MCU-Link, donde lo recibe la GUI de MATLAB. Las teclas y el touchpad sólo
modifican la configuración: SW3 la frecuencia de muestreo (match del
CTIMER0 y color del LED), SW2 el modo y el touchpad el tamaño $N$ del
próximo bloque.

\begin{figure}[H]
\centering
\resizebox{\linewidth}{!}{%
\begin{tikzpicture}[x=1cm,y=1cm, every node/.append style={font=\footnotesize}]
  \tikzset{bloque/.append style={minimum width=2.3cm}}
  % --- Fila 1: adquisicion ---
  \node[externo] (gen)  at (0,0)      {Generador\\de señal};
  \node[hw]      (adc)  at (3.4,0)    {\periferico{LPADC0}\\canal A0, \qty{16}{\bit}};
  \node[sw]      (isr)  at (6.8,0)    {ISR\\del ADC0};
  \node[mem]     (bin)  at (10.2,0)   {\texttt{buffer\_in}\\$N \times$ q15};
%
  % --- Fila 2 ---
  \node[hw]      (mux)  at (3.4,-2.4)  {\periferico{INPUTMUX}};
  \node[sw]      (proc) at (10.2,-2.4) {\texttt{procesar\_buffer()}\\FFT o \emph{bypass}};
  \node[mem]     (tab)  at (13.8,-2.4) {Tablas CMSIS\\$W_N^k$ y reordenamiento\\(flash)};
%
  % --- Fila 3 ---
  \node[externo] (p19)  at (0,-4.8)    {\pin{P0\_19}\\punto de prueba};
  \node[hw]      (tmr)  at (3.4,-4.8)  {\periferico{CTIMER0}\\match 3};
  \node[sw]      (isk)  at (6.8,-4.8)  {ISR de\\teclas};
  \node[mem]     (bout) at (10.2,-4.8) {\texttt{buffer\_out}\\$2N \times$ float32};
  \node[sw]      (lazo) at (13.8,-4.8) {Lazo principal\\(\texttt{main()})};
%
  % --- Fila 4 ---
  \node[hw]      (led)  at (3.4,-7.2)  {LED RGB};
  \node[hw]      (sw)   at (6.8,-7.2)  {Teclas SW3 ($f_s$)\\y SW2 (modo)};
  \node[hw]      (uart) at (10.2,-7.2) {\periferico{LP\_FLEXCOMM4}\\LPUART, \num{115200}~bit/s};
  \node[hw]      (tsi)  at (13.8,-7.2) {Touchpad E1\\\periferico{TSI0}};
%
  % --- Fila 5 ---
  \node[externo] (pc)   at (10.2,-9.6) {PC\\GUI de MATLAB};
%
  % --- Camino de los datos ---
  \draw[senal] (gen)  -- node[etiqueta,above]{$x(t)$} (adc);
  \draw[senal] (adc)  -- node[etiqueta,above]{FIFO0} (isr);
  \draw[senal] (isr)  -- node[etiqueta,above]{$x[n]$} (bin);
  \draw[senal] (bin)  -- node[etiqueta,right]{\texttt{buffer\_listo}} (proc);
  \draw[senal] (proc) -- node[etiqueta,right]{$X[k]$} (bout);
  \draw[senal] (bout) -- node[etiqueta,right]{trama ASCII} (uart);
  \draw[senal] (uart) -- node[etiqueta,right]{USB (COM virtual)} (pc);
%
  % --- Disparo y tablas ---
  \draw[flecha] (tab) -- (proc);
  \draw[flecha] (tmr) -- node[etiqueta,right]{MAT3 (toggle)} (mux);
  \draw[flecha] (mux) -- node[etiqueta,right]{trigger 0} (adc);
  \draw[flecha] (tmr) -- (p19);
%
  % --- Control ---
  \draw[control] (sw)   -- node[etiqueta,right]{flanco $\downarrow$} (isk);
  \draw[control] (isk)  -- node[etiqueta,above]{nuevo MR} (tmr);
  \draw[control] (isk)  -- node[etiqueta,above,sloped]{color} (led);
  \draw[control] (isk)  -- node[etiqueta,right,pos=0.3]{\texttt{cambio\_fs},\\modo} (isr);
  \draw[control] (tsi)  -- node[etiqueta,right]{toque} (lazo);
  \draw[control] (lazo) -- node[etiqueta,above,sloped]{llama} (proc);
  \draw[control] (lazo) -- node[etiqueta,below,sloped]{envía} (uart);
  \draw[control] (lazo.east) -- ++(0.7,0) |- node[etiqueta,above,pos=0.75]{$N$
                 (\texttt{BUFFER\_SIZE})} (bin.east);
%
  % --- Leyenda ---
  \begin{scope}[shift={(1.6,-11.0)},every node/.style={font=\scriptsize}]
    \foreach \c/\t [count=\k from 0] in {bloquehw/Periférico,bloquesw/Firmware,
                                          bloquemem/Memoria}
      {\fill[\c,draw=azulunc] (2.4*\k,0) rectangle ++(0.35,0.25);
       \node[right] at (2.4*\k+0.4,0.125) {\t};}
    \draw[grismedio,dashed] (7.2,0) rectangle ++(0.35,0.25);
    \node[right] at (7.6,0.125) {Externo};
  \end{scope}
\end{tikzpicture}}
\caption{Diagrama de bloques de la aplicación. Trazo continuo azul: datos;
negro: disparo y tablas; punteado: control.}
\label{fig:bloques}
\end{figure}
%
\subsection{Velocidades de muestreo}

La generación de las velocidades no cambió respecto del Laboratorio 1: con el
timer contando a \ftimer{}, reiniciándose en el match y con la salida en
toggle,
\begin{equation}
  f_s = \frac{f_{\text{timer}}}{2\,(\text{MR} + 1)}.
  \label{eq:fs}
\end{equation}
La tabla~\ref{tab:fs} repite los valores de MR (arreglo \texttt{sample\_vel})
y el color del LED de cada estado. La $f_s$ real es la que el firmware envía
en la cabecera de cada trama (sección~\ref{sec:protocolo}), así que el eje
de frecuencias de la GUI se construye con ella y no con la nominal.

\begin{table}[H]
\centering
\caption{Velocidades de muestreo, valor de match programado y color del LED.}
\label{tab:fs}
\begin{tablamuestreo}
  \filafs{8000} {9374}{Verde}   {0}{1}{0}
  \filafs{16000}{4687}{Amarillo}{1}{1}{0}
  \filafs{22000}{3408}{Azul}    {0}{0}{1}
  \filafs{44000}{1704}{Rojo}    {1}{0}{0}
  \filafs{48000}{1562}{Blanco}  {1}{1}{1}
\end{tablamuestreo}
\end{table}

\subsection{Resolución y duración de la ventana}\label{sec:resolucion}

Con las cinco velocidades y los tres tamaños hay 15 combinaciones. La
tabla~\ref{tab:resolucion} aplica la ecuación~\eqref{eq:resolucion} a cada
una, con la $f_s$ real.

\begin{table}[H]
\centering
\caption{Resolución en frecuencia $\Delta f = f_s/N$ y duración de la ventana
$T_v = N/f_s$ para cada combinación.}
\label{tab:resolucion}
\begin{tabular}{@{}l r r r r r r@{}}
\toprule
& \multicolumn{3}{c}{\textbf{$\Delta f$ [Hz]}} &
  \multicolumn{3}{c}{\textbf{$T_v$ [ms]}}\\
\cmidrule(lr){2-4}\cmidrule(l){5-7}
\textbf{$f_s$} & \textbf{512} & \textbf{1024} & \textbf{2048} &
\textbf{512} & \textbf{1024} & \textbf{2048}\\
\midrule
\filaresol{8000} {9374}
\filaresol{16000}{4687}
\filaresol{22000}{3408}
\filaresol{44000}{1704}
\filaresol{48000}{1562}
\bottomrule
\end{tabular}
\end{table}

Se observan tres cosas:
\begin{itemize}
  \item A $N$ fijo, aumentar $f_s$ amplía el rango analizado (hasta
        $f_s/2$) pero empeora la resolución en la misma proporción: a
        \qty{48}{\kilo\sps} y $N = 512$ los coeficientes están separados
        casi \qty{94}{\hertz}, seis veces más que a \qty{8}{\kilo\sps}.
  \item A $f_s$ fija, duplicar $N$ duplica la duración de la ventana y
        divide por dos $\Delta f$. También duplica la memoria, algo más que
        duplica el cálculo (tabla~\ref{tab:complejidad}) y duplica el
        tiempo de transmisión (sección~\ref{sec:tiempos}).
  \item Combinaciones distintas con el mismo cociente $f_s/N$ dan la misma
        resolución: \qty{8}{\kilo\sps} con 512 muestras y
        \qty{16}{\kilo\sps} con 1024 tienen ambas $\Delta f \approx
        \qty{15.6}{\hertz}$ y una ventana de \qty{64}{\milli\second}; la
        segunda cubre el doble de ancho de banda con el doble de
        coeficientes.
\end{itemize}

\subsection{Máquina de estados}

La configuración se modela con tres variables de estado independientes
(figura~\ref{fig:estados}):

\begin{itemize}[nosep]
  \item \texttt{state\_counter}: índice de velocidad (E0 a E4). Avanza en
        forma circular con SW3; cada cambio reprograma el match del CTIMER0,
        cambia el color del LED y descarta el bloque en curso.
  \item \texttt{modo\_seleccionado}: \emph{bypass} (M0) o FFT (M1). Alterna
        con SW2.
  \item \texttt{BUFFER\_SIZE}: tamaño de la ventana (N0 a N2: 512, 1024 y
        2048). Avanza en forma circular con cada toque del touchpad.
\end{itemize}

Las tres actúan en cualquier estado. Al reiniciar la placa el sistema
arranca en E0 (\qty{8}{\kilo\sps}), en \emph{bypass} y con $N = 512$. Ninguna
de las tres se aplica a mitad de un bloque: la rutina del ADC las copia al
comenzar cada bloque (sección~\ref{sec:bloques}).

\begin{figure}[H]
\centering
\resizebox{\linewidth}{!}{%
\begin{tikzpicture}
  % --- Anillo de velocidades ---
  \foreach \i/\f/\r/\g/\b in {0/8/0/1/0, 1/16/1/1/0, 2/22/0/0/1,
                                3/44/1/0/0, 4/48/1/1/1} {
    \node[estado] (e\i) at ({90-72*\i}:2.5)
      {\textbf{E\i}\\[-1pt]{\footnotesize\f\,kS/s}\\[1pt]\led{\r}{\g}{\b}};
  }
  \foreach \i/\j in {0/1,1/2,2/3,3/4,4/0}
    \draw[flecha] (e\i) to[bend left=18] node[etiqueta,sloped,above]{SW3} (e\j);
  \node[etiqueta] at (0,0) {avance\\circular};
  \draw[flecha] (-2.2,3.3) node[left,etiqueta]{reset} -- (e0);
%
  % --- Modo ---
  \begin{scope}[xshift=6.3cm]
    \node[estado] (m0) at (0,1.5)  {\textbf{M0}\\[-1pt]{\scriptsize\emph{bypass}}};
    \node[estado] (m1) at (0,-1.5) {\textbf{M1}\\[-1pt]{\scriptsize FFT}};
    \draw[flecha] (m0) to[bend left=35] node[etiqueta,right]{SW2} (m1);
    \draw[flecha] (m1) to[bend left=35] node[etiqueta,left]{SW2} (m0);
    \draw[flecha] (-1.4,3.0) node[left,etiqueta]{reset} -- (m0);
  \end{scope}
%
  % --- Tamano de la ventana ---
  \begin{scope}[xshift=11.4cm]
    \foreach \i/\n in {0/512, 1/1024, 2/2048} {
      \node[estado] (n\i) at ({90-120*\i}:2.0)
        {\textbf{N\i}\\[-1pt]{\footnotesize\n}};
    }
    \foreach \i/\j in {0/1,1/2,2/0}
      \draw[flecha] (n\i) to[bend left=18] node[etiqueta,sloped,above]{toque} (n\j);
    \node[etiqueta] at (0,0) {avance\\circular};
    \draw[flecha] (-2.0,2.9) node[left,etiqueta]{reset} -- (n0);
  \end{scope}
\end{tikzpicture}}
\caption{Izquierda: selección circular de la velocidad de muestreo con SW3 y
color del LED en cada estado. Centro: modo, alternado con SW2. Derecha:
tamaño de la ventana, recorrido con el touchpad E1.}
\label{fig:estados}
\end{figure}
%
\subsection{Procesamiento por bloques}\label{sec:bloques}

En el Laboratorio~2 cada muestra se procesaba en la misma interrupción en
que llegaba. La FFT, en cambio, necesita las $N$ muestras de la ventana
antes de calcular cualquier coeficiente, así que el procesamiento es por
bloques y se divide entre dos contextos que se pasan el buffer con una
bandera, \texttt{buffer\_listo} (figura~\ref{fig:linea-tiempo}):

\begin{itemize}
  \item \textbf{Rutina del ADC (productor).} Mientras
        \texttt{buffer\_listo} es falso, guarda cada muestra en
        \texttt{buffer\_in[i]} e incrementa $i$. Al llegar a $N$ pone $i$
        en cero y \texttt{buffer\_listo} en 1. Mientras
        \texttt{buffer\_listo} es verdadero, sigue leyendo la FIFO (para que
        no se desborde) pero descarta las muestras: el bloque no se toca.
  \item \textbf{Lazo principal (consumidor).} Cuando ve
        \texttt{buffer\_listo} en 1, procesa y transmite el bloque, y recién
        al terminar la pone en 0, lo que habilita la adquisición del
        siguiente.
\end{itemize}

Como sólo un contexto escribe el buffer en cada momento, no hace falta
duplicarlo ni deshabilitar interrupciones. El costo es que la adquisición
queda en pausa durante el procesamiento y la transmisión: cada espectro es
una \emph{instantánea} de $T_v$ segundos, y entre dos instantáneas se pierde
el tiempo de transmisión.

\begin{figure}[H]
\centering
\resizebox{\linewidth}{!}{%
\begin{tikzpicture}[x=1cm,y=1cm, every node/.append style={font=\footnotesize}]
  % --- Rotulos de las filas ---
  \node[etiqueta,left] at (-0.1,2.2) {ISR del ADC};
  \node[etiqueta,left] at (-0.1,1.1) {Lazo principal};
  \node[etiqueta,left] at (-0.1,0.0) {\texttt{buffer\_listo}};
%
  % --- ISR del ADC ---
  \draw[fill=bloquehw,draw=azulunc] (0,1.9) rectangle (2.6,2.5);
  \node[etiqueta] at (1.3,2.2) {Bloque $k$: $N/f_s$};
  \draw[fill=grisclaro,draw=grismedio] (2.6,1.9) rectangle (11.6,2.5);
  \node[etiqueta] at (7.1,2.2) {Lee la FIFO y descarta las muestras};
  \draw[fill=bloquehw,draw=azulunc] (11.6,1.9) rectangle (14.2,2.5);
  \node[etiqueta] at (12.9,2.2) {Bloque $k+1$};
%
  % --- Lazo principal ---
  \draw[fill=bloquesw,draw=azulunc] (2.6,0.8) rectangle (3.5,1.4);
  \node[etiqueta] at (3.05,1.1) {FFT};
  \draw[fill=bloquemem,draw=azulunc] (3.5,0.8) rectangle (11.6,1.4);
  \node[etiqueta] at (7.55,1.1) {Transmisión por la UART: cabecera, $N$ filas y \texttt{END}};
%
  % --- Bandera ---
  \draw[line width=0.9pt,draw=azulunc]
       (0,-0.25) -- (2.6,-0.25) -- (2.6,0.25) -- (11.6,0.25) -- (11.6,-0.25)
       -- (14.2,-0.25);
%
  % --- Eje de tiempo ---
  \draw[flecha] (0,-0.8) -- (14.6,-0.8) node[right,etiqueta]{$t$};
  \foreach \x in {2.6,11.6}
    \draw[grismedio,dashed] (\x,-0.8) -- (\x,2.6);
  \node[etiqueta,below] at (2.6,-0.8)  {fin del bloque $k$};
  \node[etiqueta,below] at (11.6,-0.8) {fin de la transmisión};
\end{tikzpicture}}
\caption{Línea de tiempo de un bloque (sin escala: la transmisión dura entre
16 y casi 100 veces más que la adquisición, tabla~\ref{tab:tiempos}).}
\label{fig:linea-tiempo}
\end{figure}
%
\paragraph{Configuración de cada bloque.} El tamaño, el modo y la $f_s$ se
pueden cambiar en cualquier momento, incluso mientras se transmite. Para
que el bloque sea coherente, la rutina del ADC copia
\texttt{state\_counter}, \texttt{modo\_seleccionado} y
\texttt{BUFFER\_SIZE} en \texttt{fs\_bloque}, \texttt{modo\_bloque} y
\texttt{buffer\_size\_bloque} al guardar la primera muestra ($i = 0$), y
desde ahí usa sólo las copias. El lazo principal también usa sólo las
copias para procesar y armar la cabecera, así que la FFT, la cantidad de
filas y la $f_s$ informada corresponden siempre a las muestras del bloque,
aunque el usuario haya cambiado la selección en el medio. Un cambio de modo
o de tamaño se ve recién en el bloque siguiente.

\paragraph{Muestras consecutivas.} La FFT supone que las $N$ muestras son
consecutivas y equiespaciadas a una única $f_s$. Dos situaciones romperían
esa hipótesis, y en las dos la rutina del ADC descarta el bloque parcial
(pone $i$ en cero) y vacía la FIFO:
\begin{itemize}[nosep]
  \item \textbf{Cambio de $f_s$}: la rutina de SW3 pone \texttt{cambio\_fs}
        en 1. Sin el descarte, un bloque podría tener una parte muestreada a
        una velocidad y otra a otra, y su espectro mezclaría dos escalas de
        frecuencia.
  \item \textbf{Pérdida de muestras}: si el LPADC señala desborde de la FIFO
        (\texttt{kLPADC\_ResultFIFO0}\allowbreak\texttt{OverflowFlag}) o un disparo que no pudo
        atender (\texttt{kLPADC\_Trigger}\allowbreak\texttt{ExceptionFlag}), faltan muestras y el
        bloque ya no es una ventana continua.
\end{itemize}
\texttt{cambio\_fs} arranca en 1, así que la primera interrupción también
vacía la FIFO de cualquier conversión previa a la inicialización.

\subsection{Formato de los datos}\label{sec:formato}

El LPADC0 trabaja en \qty{16}{\bit}, igual que en el Laboratorio~2, y la
rutina del ADC centra el resultado restando la mitad de escala:
\begin{equation}
  x_{q15}[n] = v - 32768 \in [-32768,\; 32767],
\end{equation}
que es directamente un \texttt{q15\_t}. El lazo principal lo pasa a punto
flotante normalizado al copiarlo a \texttt{buffer\_out}, con la parte
imaginaria en cero:
\begin{equation}
  \mathrm{Re}\,\{x[n]\} = \frac{x_{q15}[n]}{32768} \in [-1,\; 1),
  \qquad \mathrm{Im}\,\{x[n]\} = 0.
\end{equation}
Así la amplitud 1 corresponde a la mitad del rango del ADC, y la amplitud de
la ecuación~\eqref{eq:amplitud} se lee en las mismas unidades.

Se eligió punto flotante en lugar de \texttt{arm\_cfft\_q15()} porque el
Cortex-M33 del MCXN947 tiene FPU de simple precisión (el proyecto compila
con \texttt{-mfpu=fpv5-sp-d16}, ABI \emph{hard}) y porque la FFT en punto
fijo tiene que escalar: como cada etapa puede duplicar la magnitud,
\texttt{arm\_cfft\_q15()} divide por 2 en las etapas para no desbordar y
entrega el resultado dividido por $N$, perdiendo $\log_2 N$ bits (11 bits de
15 para $N = 2048$). Con \texttt{float32\_t}, de 24 bits de mantisa y
exponente propio, el resultado sin normalizar (hasta $N/2 = 1024$ veces la
amplitud de entrada) se representa sin pérdida apreciable y el piso de ruido
queda determinado por el ADC, no por la aritmética.

\subsection{Protocolo de comunicación}\label{sec:protocolo}

Cada bloque se envía como una trama de texto de $N+2$ líneas terminadas en
\verb|\r\n| (tabla~\ref{tab:trama}).

\begin{table}[H]
\centering
\caption{Formato de una trama.}
\label{tab:trama}
\small
\begin{tabular}{@{}l l p{6.2cm}@{}}
\toprule
\textbf{Línea} & \textbf{Formato} & \textbf{Contenido}\\
\midrule
Cabecera & \texttt{TP3,<seq>,<modo>,<N>,<fs\_mHz>}
         & Número de trama, modo (0: \emph{bypass}, 1: FFT), tamaño $N$ y
           $f_s$ real del bloque en milihertz\\
Filas $k = 0, \dots, N-1$ & \texttt{<k>,<RE>,<IM>}
         & Índice y bits IEEE-754 de la parte real e imaginaria, en
           hexadecimal de 8 dígitos\\
Fin      & \texttt{END} & Fin de la trama\\
\bottomrule
\end{tabular}
\end{table}

Por ejemplo, una trama FFT de 512 puntos a \qty{8}{\kilo\sps} empieza y
termina así (los valores de las filas son ilustrativos):

\begin{lstlisting}[numbers=none,caption={Ejemplo de trama.},label={lst:trama}]
TP3,17,1,512,8000000
0,3C4985F0,00000000
1,BE800000,3F000000
...
64,43800000,C1480000
...
511,BE800000,BF000000
END
\end{lstlisting}

Las decisiones del formato son:
\begin{itemize}
  \item \textbf{La $f_s$ viaja en la cabecera}, calculada con la
        ecuación~\eqref{eq:fs} a partir del MR del bloque y en milihertz,
        para que sea un entero exacto con tres decimales
        (\num{15998293}~mHz en E1). El cálculo se hace en 64 bits, porque
        $1000 \cdot f_{\text{timer}} = \num{1.5e11}$ no entra en 32. La GUI
        no necesita conocer la tabla de velocidades.
  \item \textbf{Los valores se envían como los bits del \texttt{float} en
        hexadecimal}, no en decimal. El proyecto usa la biblioteca Redlib
        con \texttt{printf} sólo de enteros (\texttt{CR\_INTEGER\_PRINTF},
        \texttt{PRINTF\_FLOAT\_ENABLE=0}), así que \texttt{snprintf} no
        puede formatear flotantes; copiar los bits a un \texttt{uint32\_t}
        con \texttt{memcpy} e imprimirlo con \verb|%08lX| no necesita
        soporte de punto flotante, transmite el valor exacto sin redondeo y
        da campos de ancho fijo, fáciles de validar. En MATLAB se recupera
        con \texttt{typecast(uint32(hex2dec(...)),'single')}.
  \item \textbf{Cada fila lleva su índice}, y la trama tiene cabecera y
        cierre. La GUI verifica que los índices sean consecutivos y que la
        trama termine con \texttt{END}; si una línea se pierde o llega
        corrupta, descarta la trama y se vuelve a sincronizar con la
        siguiente cabecera.
  \item \textbf{Se envía el espectro completo} ($N$ coeficientes, incluida la
        mitad conjugada), que es exactamente el contenido de
        \texttt{buffer\_out}. En \emph{bypass} el mismo formato lleva las
        $N$ muestras de la ventana, con la parte imaginaria en cero.
\end{itemize}

La UART se escribe con \texttt{LPUART\_WriteBlocking()}, que espera a que
haya lugar en la FIFO de transmisión. Como esa espera ocurre en el lazo
principal y no en una interrupción, no afecta a la adquisición.

\subsection{Restricción de tiempo}\label{sec:tiempos}

A diferencia del Laboratorio 2, aquí no hay que terminar el cálculo antes
de la muestra siguiente: el bloque está protegido por
\texttt{buffer\_listo} y el lazo principal puede tardar lo que necesite. Lo
que se ve afectado es el período con que se refresca el espectro,
\begin{equation}
  T_{\text{bloque}} = T_v + T_{\text{FFT}} + T_{\text{tx}},
  \qquad
  T_{\text{tx}} = \frac{\text{bytes de la trama} \cdot 10}{\num{115200}},
  \label{eq:tbloque}
\end{equation}
donde cada byte ocupa 10 bits en la línea (inicio, 8 datos y parada). Una
fila tiene entre 21 y 24 bytes según la cantidad de dígitos del índice. La
tabla~\ref{tab:tiempos} aplica la ecuación~\eqref{eq:tbloque} a la menor y a
la mayor $f_s$, sin $T_{\text{FFT}}$.

\begin{table}[H]
\centering
\caption{Bytes por trama, tiempo de transmisión a \num{115200}~bit/s, duración
de la ventana y período de refresco del espectro.}
\label{tab:tiempos}
\begin{tabular}{@{}r r r r r r r@{}}
\toprule
& & & \multicolumn{2}{c}{\textbf{$T_v$ [ms]}} &
      \multicolumn{2}{c}{\textbf{$T_{\text{bloque}}$ [s]}}\\
\cmidrule(lr){4-5}\cmidrule(l){6-7}
\textbf{$N$} & \textbf{Bytes} & \textbf{$T_{\text{tx}}$ [s]} &
\textbf{\qty{8}{\kilo\sps}} & \textbf{\qty{48}{\kilo\sps}} &
\textbf{\qty{8}{\kilo\sps}} & \textbf{\qty{48}{\kilo\sps}}\\
\midrule
\filatx{512} {11693}
\filatx{1024}{23494}
\filatx{2048}{48070}
\bottomrule
\end{tabular}
\end{table}

La transmisión domina: dura entre 16 y casi 100 veces la ventana, y el espectro se
refresca aproximadamente una vez por segundo con $N = 512$ y una vez cada
cuatro segundos con $N = 2048$, casi sin importar la $f_s$.

El cálculo de la FFT es mucho más corto. Con la estimación habitual de
$5\,N\log_2 N$ operaciones de punto flotante (cada mariposa radix-2: un
producto complejo, cuatro productos y dos sumas, más dos sumas complejas),
la FFT de \num{2048} puntos son unas \num{112640} operaciones, del orden de
\qty{1}{\milli\second} a \qty{150}{\mega\hertz} aun contando los accesos a
memoria y la compilación sin optimizar.

\paragraph{¿Se podría transmitir en forma continua?} Para no perder
muestras, cada trama tendría que enviarse en menos tiempo que $T_v$, es
decir, a una tasa de $(\text{bytes por coeficiente}) \cdot f_s$. Con el
formato actual (unos \num{23.5} bytes por coeficiente) eso es
\qty{1.9}{\mega\bit\per\second} a \qty{8}{\kilo\sps} y
\qty{11.3}{\mega\bit\per\second} a \qty{48}{\kilo\sps}. Incluso con el
formato más compacto posible, en binario y sólo con los $N/2$ coeficientes
no redundantes (8 bytes cada uno, 4 bytes por muestra de entrada), harían
falta \qty{320}{\kilo\bit\per\second} y \qty{1.92}{\mega\bit\per\second}.
Las dos cifras superan los \num{115200}~bit/s del enlace, así que el
esquema de instantáneas no es una simplificación sino la única opción con
este enlace. Para un análisis continuo habría que reducir lo que se
transmite (por ejemplo, sólo el módulo de unos pocos coeficientes) o
promediar varios espectros en la placa antes de enviar uno.

\subsection{Uso de memoria}

\texttt{buffer\_in} y \texttt{buffer\_out} se dimensionan una sola vez para
el máximo, \texttt{BUFFER\_MAX\_SIZE} = \num{2048}, y los bloques
más chicos usan sólo su parte inicial. \texttt{buffer\_in} ocupa
4~KiB ($2048 \times 2$ bytes) y \texttt{buffer\_out}
16~KiB ($2 \times 2048 \times 4$ bytes), unos
20~KiB de RAM en total. La FFT trabaja en el lugar sobre
\texttt{buffer\_out}, sin buffers auxiliares, y las tablas de la
tabla~\ref{tab:tablas} están en flash. Que la entrada y la salida sean
buffers distintos es lo que pide la consigna, y además permite conservar
las muestras originales en \texttt{buffer\_in} aunque la FFT sobrescriba
\texttt{buffer\_out}.

%%%%%%%%%%%%%%%%%%%%%%%%%%%%%%%%%%%%%%%%%%%%%%%%%%%%%%%%%%%%%%%%%%%%%%%%%%%%%%%
\ejercicio{Implementación}\label{sec:implementacion}
%%%%%%%%%%%%%%%%%%%%%%%%%%%%%%%%%%%%%%%%%%%%%%%%%%%%%%%%%%%%%%%%%%%%%%%%%%%%%%%

La lógica propia está en \archivo{source/TP3.c}; el resto del proyecto es
código generado por Config Tools, drivers del SDK y CMSIS-DSP. Los scripts
de MATLAB están en \archivo{TP3/matlab/}. El código completo está en el
Anexo~\ref{sec:codigo}.

\subsection{Constantes y variables globales}

Los valores de match se mantienen en \texttt{sample\_vel}, igual que en los
laboratorios anteriores. \texttt{BUFFER\_MAX\_SIZE} (\num{2048}) fija la
capacidad de los buffers. \texttt{BUFFER\_SIZE} conserva el nombre de la
constante del Laboratorio~2 (\texttt{\#define BUFFER\_SIZE 512}), pero ahora es una variable \texttt{volatile}
que modifica el touchpad. Las variables que comparten las interrupciones y
el lazo principal (\texttt{state\_counter}, \texttt{modo\_seleccionado},
\texttt{BUFFER\_SIZE}, \texttt{buffer\_listo} y \texttt{cambio\_fs}) se
declaran \texttt{volatile}. \texttt{fs\_bloque}, \texttt{modo\_bloque} y
\texttt{buffer\_size\_bloque} son la configuración del bloque
(sección~\ref{sec:bloques}): la escribe la rutina del ADC al comenzar un
bloque y la lee el lazo principal sólo mientras \texttt{buffer\_listo} es
verdadero. \texttt{pte\_fft} apunta a la instancia de CMSIS-DSP del tamaño
activo. El estado del touchpad (\texttt{touch\_base},
\texttt{touch\_tocado} y los demás \texttt{touch\_*}) y la base de tiempo
\texttt{ms\_ticks} son los del Laboratorio~2.

\subsection{Programa principal}

\texttt{main()} inicializa pines y relojes, enciende el VREF
(sección~\ref{sec:cambios-perif}) y recién después inicializa los
periféricos generados por Config Tools (CTIMER0, LPADC0, GPIO0 y la UART).
A continuación fija el color del LED para E0, configura el SysTick para
interrumpir cada \qty{1}{\milli\second}, mide la base del touchpad
(\texttt{init\_touch()}), envía la línea \texttt{TP3 boot} (que la GUI
ignora) y entra en el lazo principal (figura~\ref{fig:flujo-lazo}).

\begin{figure}[H]
\centering
\begin{tikzpicture}[node distance=0.55cm and 0.4cm]
  \coordinate (izq) at (-3.8,0);
  \coordinate (der) at (3.0,0);
  \node[inicio] (m0) at (0,0) {\texttt{main()}};
  \node[proceso, below=of m0] (m1)
     {Pines, relojes, VREF y periféricos;\\SysTick de \qty{1}{\milli\second};
      \texttt{init\_touch()}};
  \node[proceso, below=of m1] (m2) {\texttt{touch\_actualizar()}};
  \node[decision, below=0.6cm of m2] (m3) {¿\texttt{buffer\_listo}?};
  \node[proceso, below=0.6cm of m3] (m4)
     {\texttt{procesar\_buffer()}:\\q15 $\rightarrow$ float32;\\FFT si el bloque es FFT};
  \node[proceso, below=of m4] (m5)
     {\texttt{enviar\_buffer()}:\\cabecera, $N$ filas y \texttt{END}\\
      (atiende el touchpad entre filas)};
  \node[proceso, below=of m5] (m6) {\texttt{buffer\_listo} $\leftarrow 0$};
  \draw[flecha] (m0) -- (m1);
  \draw[flecha] (m1) -- (m2);
  \draw[flecha] (m2) -- (m3);
  \draw[flecha] (m3) -- node[etiqueta,right]{sí} (m4);
  \draw[flecha] (m4) -- (m5);
  \draw[flecha] (m5) -- (m6);
  \draw[flecha] (m6.west) -- (m6.west -| izq) |- (m2.west);
  \draw[flecha] (m3.east) -- node[etiqueta,above]{no} (m3.east -| der) |- (m2.east);
\end{tikzpicture}
\caption{Lazo principal: touchpad, procesamiento y transmisión.}
\label{fig:flujo-lazo}
\end{figure}

\subsection{Rutina del ADC}

El código~\ref{lst:isr} muestra el núcleo de la rutina del ADC y la
figura~\ref{fig:flujo-adc} su diagrama de flujo. Primero se atiende el
descarte por cambio de $f_s$ o pérdida de muestras; si no corresponde, se
lee la conversión, se centra a q15 y, si el bloque no está protegido, se
guarda.

\begin{lstlisting}[caption={Adquisición de un bloque en \texttt{ADC0\_IRQHANDLER()} (\texttt{source/TP3.c}).},label={lst:isr}]
if (cambio_fs ||
    (status_flag & (kLPADC_ResultFIFO0OverflowFlag |
                    kLPADC_TriggerExceptionFlag))) {
    index_buffer = 0;                 /* descarta el bloque parcial */
    cambio_fs = false;
    while (LPADC_GetConvResult(ADC0_PERIPHERAL, &result, 0)) {
    }                                 /* vacia la FIFO */

} else if (LPADC_GetConvResult(ADC0_PERIPHERAL, &result, 0)) {
    q15_t muestra = (q15_t)((int32_t)result.convValue - 32768);

    if (!buffer_listo) {
        if (index_buffer == 0) {      /* configuracion de ESTE bloque */
            fs_bloque = state_counter;
            modo_bloque = modo_seleccionado;
            buffer_size_bloque = BUFFER_SIZE;
        }
        buffer_in[index_buffer] = muestra;
        index_buffer++;
        if (index_buffer == buffer_size_bloque) {
            index_buffer = 0;
            buffer_listo = true;      /* entrega el bloque a main */
        }
    }
}
\end{lstlisting}

\begin{figure}[H]
\centering
\begin{tikzpicture}[node distance=0.5cm and 0.4cm]
  \coordinate (busder) at (7.6,0);
  \coordinate (busizq) at (-3.9,0);
  \node[inicio] (a0) at (0,0) {\texttt{ADC0\_IRQHandler}};
  \node[proceso, below=of a0] (a1) {Leer y limpiar banderas\\de trigger y de estado};
  \node[decision, below=of a1] (a2) {¿Cambio de $f_s$\\o desborde?};
  \node[proceso] (a3) at ([xshift=5cm]a2) {$i \leftarrow 0$; \texttt{cambio\_fs} $\leftarrow 0$\\vaciar la FIFO0};
  \node[decision, below=0.6cm of a2] (a4) {¿Hay resultado\\en la FIFO0?};
  \node[proceso, below=0.6cm of a4] (a5) {$x \leftarrow \mathrm{res} - 32768$};
  \node[decision, below=0.6cm of a5] (a6) {¿\texttt{buffer\_listo}?};
  \node[decision, below=0.6cm of a6] (a7) {¿$i = 0$?};
  \node[proceso] (a8) at ([xshift=5cm]a7) {Copiar $f_s$, modo y $N$\\del bloque};
  \node[proceso, below=0.6cm of a7] (a9) {\texttt{buffer\_in}$[i] \leftarrow x$\\$i \leftarrow i + 1$};
  \node[decision, below=0.6cm of a9] (a10) {¿$i = N_{\text{bloque}}$?};
  \node[proceso] (a11) at ([xshift=5cm]a10) {$i \leftarrow 0$\\\texttt{buffer\_listo} $\leftarrow 1$};
  \node[inicio, below=0.7cm of a10] (a12) {Retorno};

  \draw[flecha] (a0) -- (a1);
  \draw[flecha] (a1) -- (a2);
  \draw[flecha] (a2) -- node[etiqueta,above]{sí} (a3);
  \draw[flecha] (a2) -- node[etiqueta,right]{no} (a4);
  \draw[flecha] (a4) -- node[etiqueta,right]{sí} (a5);
  \draw[flecha] (a5) -- (a6);
  \draw[flecha] (a6) -- node[etiqueta,right]{no} (a7);
  \draw[flecha] (a7) -- node[etiqueta,above]{sí} (a8);
  \draw[flecha] (a8) |- (a9);
  \draw[flecha] (a7) -- node[etiqueta,right]{no} (a9);
  \draw[flecha] (a9) -- (a10);
  \draw[flecha] (a10) -- node[etiqueta,above]{sí} (a11);
  \draw[flecha] (a10) -- node[etiqueta,right]{no} (a12);
  % Salidas hacia el retorno por la derecha
  \draw[flecha] (a3.east) -- (a3.east -| busder) |- (a12.east);
  \draw (a11.east) -- (a11.east -| busder);
  % Salidas hacia el retorno por la izquierda
  \draw[flecha] (a4.west) -- node[etiqueta,above]{no} (a4.west -| busizq) |- (a12.west);
  \draw (a6.west) -- node[etiqueta,above]{sí} (a6.west -| busizq);
\end{tikzpicture}
\caption{Diagrama de flujo de la rutina de fin de conversión del ADC. Con
\texttt{buffer\_listo} en 1 la muestra se lee (la FIFO no se llena) pero se
descarta.}
\label{fig:flujo-adc}
\end{figure}

\subsection{Procesamiento del bloque}

\texttt{procesar\_buffer()} (código~\ref{lst:proc}) convierte el bloque a
complejo en punto flotante y, si el bloque se adquirió en modo FFT, elige
la instancia de CMSIS-DSP del tamaño del bloque y calcula la transformada
en el lugar. \texttt{init\_fft()} conserva el nombre de la versión inicial,
en la que el tamaño era una constante de compilación y la instancia se
elegía una sola vez con \texttt{\#if}; ahora se llama en cada bloque y elige
con \texttt{buffer\_size\_bloque}, el tamaño con que se adquirió, y no con
\texttt{BUFFER\_SIZE}, que el touchpad puede haber cambiado mientras tanto.
Elegir la instancia no cuesta nada: es asignar un puntero a una estructura
constante.

\begin{lstlisting}[caption={Selección de la instancia y cálculo de la FFT (\texttt{source/TP3.c}).},label={lst:proc}]
void init_fft(void) {
    switch (buffer_size_bloque) {
        case 512U:  pte_fft = &arm_cfft_sR_f32_len512;  break;
        case 1024U: pte_fft = &arm_cfft_sR_f32_len1024; break;
        case 2048U: pte_fft = &arm_cfft_sR_f32_len2048; break;
    }
}

void procesar_buffer(void) {
    for (uint16_t i = 0; i < buffer_size_bloque; i++) {
        buffer_out[2*i]   = (float32_t)buffer_in[i] / 32768.0f;  /* Re */
        buffer_out[2*i+1] = 0.0f;                                 /* Im */
    }
    if (modo_bloque == MODO_FFT) {
        init_fft();
        arm_cfft_f32(pte_fft, buffer_out, 0, 1);  /* directa, orden natural */
    }
}
\end{lstlisting}

En \emph{bypass} la función sólo hace la copia: el buffer de salida es
distinto del de entrada, como pide la consigna, y contiene las muestras sin
transformar.

\subsection{Transmisión}

\texttt{enviar\_buffer()} (código~\ref{lst:tx}) arma la cabecera con la
configuración del bloque y envía una fila por coeficiente. Entre fila y
fila llama a \texttt{touch\_actualizar()}: la transmisión de un bloque de
\num{2048} puntos tarda más de cuatro segundos
(tabla~\ref{tab:tiempos}), y sin esa llamada el touchpad no respondería
durante ese tiempo. Como el touchpad sólo cambia \texttt{BUFFER\_SIZE}, la
trama en curso no se altera.

\begin{lstlisting}[caption={Armado y envío de la trama en \texttt{enviar\_buffer()} (\texttt{source/TP3.c}).},label={lst:tx}]
uint32_t fs_millihz =
    (uint32_t)((1000ULL * CTIMER0_TICK_FREQ) /
    (2ULL * (CTIMER0_config.prescale + 1ULL) *
     (sample_vel[fs_bloque] + 1ULL)));

snprintf(linea, sizeof(linea), "TP3,%lu,%u,%u,%lu\r\n",
         (unsigned long)++numero_bloque, (unsigned)modo_bloque,
         (unsigned)buffer_size_bloque, (unsigned long)fs_millihz);
if (!enviar_texto_uart(linea)) return;

for (uint16_t i = 0; i < buffer_size_bloque; i++) {
    uint32_t real, imag;
    memcpy(&real, &buffer_out[2*i],   sizeof(real));  /* bits del float */
    memcpy(&imag, &buffer_out[2*i+1], sizeof(imag));
    snprintf(linea, sizeof(linea), "%u,%08lX,%08lX\r\n", (unsigned)i,
             (unsigned long)real, (unsigned long)imag);
    if (!enviar_texto_uart(linea)) return;
    touch_actualizar(ms_ticks);     /* el touchpad no espera a END */
}
(void)enviar_texto_uart("END\r\n");
\end{lstlisting}

Se usa \texttt{memcpy} en lugar de un \emph{cast} de punteros
(\texttt{*(uint32\_t *)\&f}) porque reinterpretar un \texttt{float} a
través de un puntero a entero viola las reglas de \emph{aliasing} de C y
el compilador podría optimizarlo de forma incorrecta; \texttt{memcpy} es la
forma portable y, para 4 bytes, el compilador la reduce a una sola
transferencia.

\subsection{Rutinas de las teclas}

La rutina de SW3 avanza \texttt{state\_counter} en forma circular,
reprograma el CTIMER0 con \texttt{setup\_new\_match()} (detiene el timer,
carga el nuevo match, lo reinicia y lo arranca, como en el Laboratorio~1),
actualiza el color del LED con \texttt{set\_state()} y pone
\texttt{cambio\_fs} en 1. La rutina de SW2 sólo alterna
\texttt{modo\_seleccionado}. Ninguna toca los buffers: el efecto se aplica
en la rutina del ADC al comenzar el bloque siguiente
(figura~\ref{fig:flujo-teclas}).

\begin{figure}[H]
\centering
\begin{tikzpicture}[node distance=0.55cm and 0.4cm]
  % ---------------- ISR de SW3 ----------------
  \node[inicio] (b0) {\texttt{GPIO0\_INT\_0} (SW3)};
  \node[proceso, below=of b0] (b1) {\texttt{state\_counter} $\leftarrow$\\(\texttt{state\_counter}+1) mod 5};
  \node[proceso, below=of b1] (b2) {\texttt{setup\_new\_match}\\(\texttt{sample\_vel[state]})};
  \node[proceso, below=of b2] (b3) {\texttt{set\_state()}: color del LED};
  \node[proceso, below=of b3] (b4) {\texttt{cambio\_fs} $\leftarrow$ 1};
  \node[proceso, below=of b4] (b5) {Limpiar bandera};
  \draw[flecha] (b0) -- (b1);
  \draw[flecha] (b1) -- (b2);
  \draw[flecha] (b2) -- (b3);
  \draw[flecha] (b3) -- (b4);
  \draw[flecha] (b4) -- (b5);

  % ---------------- ISR de SW2 ----------------
  \begin{scope}[xshift=6.5cm]
    \node[inicio] (c0) {\texttt{GPIO0\_INT\_1} (SW2)};
    \node[proceso, below=of c0] (c1) {\texttt{modo\_seleccionado} $\leftarrow$\\FFT si era \emph{bypass},\\\emph{bypass} si era FFT};
    \node[proceso, below=of c1] (c2) {Limpiar bandera};
    \draw[flecha] (c0) -- (c1);
    \draw[flecha] (c1) -- (c2);
  \end{scope}
\end{tikzpicture}
\caption{Diagramas de flujo de las rutinas de las teclas. Izquierda: cambio
de velocidad (SW3). Derecha: cambio de modo (SW2).}
\label{fig:flujo-teclas}
\end{figure}

\subsection{Touchpad}

El touchpad se lee exactamente como en el Laboratorio~2, desde el lazo
principal y sin interrupciones: al arrancar, \texttt{init\_touch()} configura
el TSI0 en modo \emph{self-cap} con disparo por software y mide la base como
el promedio de 16 barridos sin tocar (por eso no hay que tocar el touchpad
durante el reset). Después, \texttt{touch\_actualizar()} dispara un barrido
cada \qty{10}{\milli\second} y consulta si terminó en las vueltas
siguientes, sin esperar; si un barrido no termina en
\qty{50}{\milli\second}, el touchpad se desactiva y el resto del programa
sigue funcionando. Una lectura cuenta como toque si está más de 20 cuentas
por debajo de la base y como liberación si vuelve a quedar a menos de 10
(histéresis), y el cambio se acepta con 3 lecturas seguidas (antirrebote de
unos \qty{30}{\milli\second}). Sólo el paso de \emph{sin tocar} a
\emph{tocado} cuenta como toque, así que mantener el dedo apoyado no avanza
varios tamaños.

Lo único que cambió es la acción: cada toque nuevo avanza
\texttt{BUFFER\_SIZE} en forma circular, 512 $\rightarrow$ 1024
$\rightarrow$ 2048 $\rightarrow$ 512. El valor se aplica al comenzar el
próximo bloque. En el TP3 se quitó el aviso por \texttt{PRINTF} cuando el
touchpad no responde, porque la UART transporta el protocolo y un texto
intercalado rompería una trama.

\subsection{Aplicación de MATLAB}\label{sec:matlab}

En \archivo{TP3/matlab/} hay tres scripts:

\begin{itemize}
  \item \archivo{recibir_tp3.m}\texttt{(puerto, baud)} es la GUI. Abre el
        puerto con \texttt{serialport} (\num{115200}~bit/s por defecto) y
        terminador \verb|CR/LF|, y en un lazo que dura mientras la ventana
        esté abierta: busca una línea que empiece con \texttt{TP3,}, valida
        la cabecera ($N \in \{512, 1024, 2048\}$, modo 0 o 1, $f_s > 0$),
        lee las $N$ filas verificando con una expresión regular el formato
        y que el índice sea el esperado, decodifica los flotantes y
        comprueba el \texttt{END}. En modo FFT dibuja el espectro de un solo
        lado de la ecuación~\eqref{eq:amplitud} entre 0 y $f_s/2$, con el eje
        en hertz construido con la $f_s$ de la cabecera; en \emph{bypass}
        dibuja las muestras contra el tiempo. El título muestra el número
        de trama, el modo, $N$ y $f_s$. Una trama con errores genera una
        advertencia y se descarta, sin cerrar la GUI.
  \item \archivo{diagnosticar_tp3.m} captura una sola trama, la guarda en
        \archivo{.mat} y \archivo{.txt}, e informa el rango de la parte real
        y la imaginaria, la cantidad de valores no nulos y no finitos, los
        primeros datos en hexadecimal y en decimal y, en modo FFT, la
        componente continua y el pico principal (frecuencia y amplitud). Se
        usó para depurar el enlace y para tomar los valores numéricos de la
        sección~\ref{sec:verificacion}.
  \item \archivo{explorar_fft_tp3.m} es una referencia en la PC: genera una
        senoidal, la cuantiza a q15 como el ADC, calcula su FFT con
        \texttt{fft} y la grafica con la misma normalización que la GUI.
        Permite prever el resultado de cada combinación de $f_s$, $N$ y
        $f_0$ (por ejemplo, la fuga de la sección~\ref{sec:fuga}).
\end{itemize}

\begin{lstlisting}[language=Matlab,caption={Decodificación de una fila y espectro de un solo lado en \texttt{recibir\_tp3.m}.},label={lst:matlab}]
tok = regexp(row,'^(\d+),([0-9A-F]{8}),([0-9A-F]{8})$','tokens','once');
if isempty(tok) || str2double(tok{1}) ~= k
    error('TP3:Frame','Trama incompleta o desordenada.');
end
re = double(typecast(uint32(hex2dec(tok{2})),'single'));
im = double(typecast(uint32(hex2dec(tok{3})),'single'));
X(k+1) = complex(re,im);
% ... al completar la trama, en modo FFT:
A = abs(X(1:N/2+1))/N;
A(2:end-1) = 2*A(2:end-1);
plot(ax,(0:N/2)*Fs/N,A);
\end{lstlisting}

%%%%%%%%%%%%%%%%%%%%%%%%%%%%%%%%%%%%%%%%%%%%%%%%%%%%%%%%%%%%%%%%%%%%%%%%%%%%%%%
\ejercicio{Verificación}\label{sec:verificacion}
%%%%%%%%%%%%%%%%%%%%%%%%%%%%%%%%%%%%%%%%%%%%%%%%%%%%%%%%%%%%%%%%%%%%%%%%%%%%%%%

\subsection{Procedimiento}

\begin{enumerate}[nosep]
  \item Conectar el generador a \texttt{J4[2]} (ADC0\_A0) con una senoidal
        de \qty{1}{\kilo\hertz} y un \emph{offset} que la mantenga dentro
        del rango del ADC. Conectar la placa a la PC por el USB del
        MCU-Link.
  \item Ejecutar \texttt{recibir\_tp3("COMx", 115200)} en MATLAB (con
        ningún otro programa usando el puerto).
  \item Con $f_s = \qty{8}{\kilo\sps}$ y $N = 512$, verificar en
        \emph{bypass} que la GUI muestra la senoidal en el tiempo; pasar a
        FFT con SW2 y verificar que el pico está en \qty{1}{\kilo\hertz}.
  \item Recorrer los tres tamaños con el touchpad a $f_s$ fija, y las cinco
        velocidades con SW3 a $N$ fijo, verificando en el título de la GUI
        que $N$ y $f_s$ cambian en el bloque siguiente.
  \item Repetir con una senoidal de \qty{5}{\kilo\hertz} a
        \qty{8}{\kilo\sps} y a \qty{16}{\kilo\sps} (aliasing) y con una
        cuadrada de \qty{500}{\hertz} (armónicos).
  \item Para los valores numéricos, cerrar la GUI y capturar una trama con
        \texttt{diagnosticar\_tp3}.
\end{enumerate}

\subsection{Verificación por simulación}

Antes de medir en la placa se calculó en la PC, con
\archivo{explorar_fft_tp3.m}, el espectro esperado para $f_s =
\qty{8}{\kilo\sps}$ y $N = 512$. Con $f_0 = \qty{1000}{\hertz}$ el script
informa una ventana de \qty{64}{\milli\second}, una separación de
\qty{15.625}{\hertz} y un único pico en \qty{1000}{\hertz} de amplitud
$0{,}5$; con $f_0 = \qty{1005}{\hertz}$ aparece la fuga de la
tabla~\ref{tab:fuga}.

\captura[0.32\textheight]{img/sim-fuga-1000.png}{MATLAB}{\archivo{explorar_fft_tp3.m}
con $f_0 = \qty{1000}{\hertz}$, $f_s = \qty{8}{\kilo\sps}$, $N = 512$: la
senoidal tiene un número entero de ciclos y cae en un solo
coeficiente.}{fig:sim-1000}

\captura[0.32\textheight]{img/sim-fuga-1005.png}{MATLAB}{\archivo{explorar_fft_tp3.m}
con $f_0 = \qty{1005}{\hertz}$: fuga espectral hacia los coeficientes
vecinos.}{fig:sim-1005}

\subsection{Transmisión y modo \emph{bypass}}

\captura[0.3\textheight]{img/terminal-trama.png}{Terminal serie}{Comienzo de
una trama recibida: cabecera, primeras filas.}{fig:terminal}

\captura{img/gui-bypass-8k.png}{GUI de MATLAB}{Modo \emph{bypass}: senoidal
de \qty{1}{\kilo\hertz} muestreada a \qty{8}{\kilo\sps}, $N = 512$
(\qty{64}{\milli\second}, 64 ciclos).}{fig:gui-bypass}

\subsection{Espectro al variar el tamaño de la ventana}

A $f_s = \qty{8}{\kilo\sps}$, \qty{1}{\kilo\hertz} es múltiplo de $\Delta f$
para los tres tamaños (coeficientes 64, 128 y 256), así que el pico debe
quedar en \qty{1000}{\hertz} con la amplitud de la senoidal.

\captura{img/gui-fft-8k-512.png}{GUI de MATLAB}{FFT de una senoidal de
\qty{1}{\kilo\hertz} a \qty{8}{\kilo\sps} con $N = 512$
($\Delta f = \qty{15.6}{\hertz}$).}{fig:gui-8k-512}

\captura{img/gui-fft-8k-1024.png}{GUI de MATLAB}{Igual que la
figura~\ref{fig:gui-8k-512} con $N = 1024$
($\Delta f = \qty{7.8}{\hertz}$).}{fig:gui-8k-1024}

\captura{img/gui-fft-8k-2048.png}{GUI de MATLAB}{Igual que la
figura~\ref{fig:gui-8k-512} con $N = 2048$
($\Delta f = \qty{3.9}{\hertz}$).}{fig:gui-8k-2048}

\subsection{Espectro al variar la frecuencia de muestreo}

A $N = 512$ y velocidades mayores, \qty{1}{\kilo\hertz} ya no es múltiplo de
$\Delta f$: el pico se desplaza al coeficiente más cercano y aparece fuga.
La tabla~\ref{tab:medido} incluye el coeficiente esperado.

\captura{img/gui-fft-22k-512.png}{GUI de MATLAB}{FFT de una senoidal de
\qty{1}{\kilo\hertz} a \qty{22}{\kilo\sps} con $N = 512$
($\Delta f = \qty{43}{\hertz}$).}{fig:gui-22k-512}

\captura{img/gui-fft-48k-512.png}{GUI de MATLAB}{FFT de una senoidal de
\qty{1}{\kilo\hertz} a \qty{48}{\kilo\sps} con $N = 512$
($\Delta f = \qty{93.7}{\hertz}$).}{fig:gui-48k-512}

\subsection{Aliasing y armónicos}

El LPADC0 no tiene filtro antialias a la entrada, así que una componente por
encima de $f_s/2$ aparece reflejada en $|f - m f_s|$. Una senoidal de
\qty{5}{\kilo\hertz} muestreada a \qty{8}{\kilo\sps} debe verse en
\qty{3}{\kilo\hertz}; a \qty{16}{\kilo\sps}, en su frecuencia real. Una
cuadrada de \qty{500}{\hertz} y amplitud $A$ tiene sólo armónicos impares,
de amplitud $4A/(\pi m)$; a \qty{8}{\kilo\sps} los que superan
\qty{4}{\kilo\hertz} se pliegan sobre la banda.

\captura{img/gui-alias-8k.png}{GUI de MATLAB}{Senoidal de
\qty{5}{\kilo\hertz} a \qty{8}{\kilo\sps}: aparece en
\qty{3}{\kilo\hertz}.}{fig:gui-alias-8k}

\captura{img/gui-alias-16k.png}{GUI de MATLAB}{La misma senoidal a
\qty{16}{\kilo\sps}: aparece en \qty{5}{\kilo\hertz}.}{fig:gui-alias-16k}

\captura{img/gui-cuadrada-8k.png}{GUI de MATLAB}{Cuadrada de
\qty{500}{\hertz} a \qty{8}{\kilo\sps}: armónicos impares y armónicos
plegados.}{fig:gui-cuadrada}

\subsection{Valores medidos}

\begin{table}[H]
\centering
\caption{Pico principal esperado y medido (con \archivo{diagnosticar_tp3.m}).
El esperado es el coeficiente más cercano a la frecuencia de entrada.}
\label{tab:medido}
\footnotesize
\begin{tabular}{@{}r r r r r r r@{}}
\toprule
\textbf{$f_s$ [Hz]} & \textbf{$N$} & \textbf{$\Delta f$ [Hz]} &
\textbf{$f_{\text{in}}$ [Hz]} & \textbf{Pico esperado [Hz]} &
\textbf{Pico medido [Hz]} & \textbf{Amplitud medida}\\
\midrule
\num{8000}     & 512  & \num{15.63} & \num{1000} & \num{1000}    & \completar{} & \completar{}\\
\num{8000}     & 1024 & \num{7.81}  & \num{1000} & \num{1000}    & \completar{} & \completar{}\\
\num{8000}     & 2048 & \num{3.91}  & \num{1000} & \num{1000}    & \completar{} & \completar{}\\
\num{22000.59} & 512  & \num{42.97} & \num{1000} & \num{988.31}  & \completar{} & \completar{}\\
\num{47984.64} & 512  & \num{93.72} & \num{1000} & \num{1030.92} & \completar{} & \completar{}\\
\num{8000}     & 512  & \num{15.63} & \num{5000} & \num{3000}    & \completar{} & \completar{}\\
\num{15998.29} & 512  & \num{31.25} & \num{5000} & \num{4999.47} & \completar{} & \completar{}\\
\bottomrule
\end{tabular}
\end{table}

\subsection{Cumplimiento de los requerimientos}

\begin{table}[H]
\centering
\caption{Verificación de los requerimientos de la consigna.}
\label{tab:cumplimiento}
\begin{tabular}{@{}l p{6.8cm} l l@{}}
\toprule
\textbf{Req.} & \textbf{Cómo se implementa} & \textbf{Evidencia} & \textbf{Resultado}\\
\midrule
\ref{req:fs}    & Cadena CTIMER0 + INPUTMUX + LPADC0 del Laboratorio~1; SW3
                  elige la $f_s$; cambio de $f_s$ con descarte del bloque
                & Figuras~\ref{fig:gui-8k-512}, \ref{fig:gui-48k-512} & \checkmark\\
\ref{req:n}     & \texttt{buffer\_in[2048]}; el touchpad elige 512, 1024 o
                  2048, aplicado al comenzar el bloque
                & Figuras~\ref{fig:gui-8k-512} a~\ref{fig:gui-8k-2048}
                & \completar{verificar}\\
\ref{req:fft}   & \texttt{arm\_cfft\_f32()} con la instancia del tamaño del
                  bloque
                & Tabla~\ref{tab:medido} & \checkmark\\
\ref{req:tecla} & SW2 alterna FFT y \emph{bypass}; el resultado va a
                  \texttt{buffer\_out}, distinto de \texttt{buffer\_in}
                & Figuras~\ref{fig:gui-bypass} y~\ref{fig:gui-8k-512} & \checkmark\\
\ref{req:sci}   & LPUART (LP\_FLEXCOMM4) a \num{115200}~bit/s, trama con
                  cabecera, $N$ filas y \texttt{END}
                & Figura~\ref{fig:terminal} & \checkmark\\
\ref{req:gui}   & \archivo{recibir_tp3.m}: espectro de un solo lado en hertz
                  o señal en el tiempo
                & Figuras~\ref{fig:gui-bypass} a~\ref{fig:gui-cuadrada} & \checkmark\\
\ref{req:concl} & Análisis de resolución, tiempos y fuga
                & Secciones~\ref{sec:resolucion}, \ref{sec:tiempos}
                  y~\ref{sec:conclusiones} & \checkmark\\
\bottomrule
\end{tabular}
\end{table}

%%%%%%%%%%%%%%%%%%%%%%%%%%%%%%%%%%%%%%%%%%%%%%%%%%%%%%%%%%%%%%%%%%%%%%%%%%%%%%%
\ejercicio{Conclusiones}\label{sec:conclusiones}
%%%%%%%%%%%%%%%%%%%%%%%%%%%%%%%%%%%%%%%%%%%%%%%%%%%%%%%%%%%%%%%%%%%%%%%%%%%%%%%

Se implementó sobre la FRDM-MCXN947 una aplicación que adquiere la señal con
la cadena del Laboratorio~1 a cinco velocidades de muestreo, la guarda en
bloques de 512, 1024 o 2048 muestras, les aplica la FFT con
\texttt{arm\_cfft\_f32()} de CMSIS-DSP (o las deja pasar sin transformar, en
\emph{bypass}) hacia un buffer de salida distinto del de entrada, y envía el
resultado por la UART a una GUI de MATLAB que dibuja el espectro o la señal.
SW3 elige la frecuencia de muestreo, SW2 el modo y el touchpad el tamaño de
la ventana, sin recompilar.

\subsection{Efecto de la frecuencia de muestreo y del tamaño de la ventana}

\begin{itemize}
  \item \textbf{La resolución depende sólo de la duración de la ventana.}
        $\Delta f = f_s/N = 1/T_v$: para distinguir dos componentes
        separadas $\delta f$ hace falta observar la señal durante al menos
        $1/\delta f$ segundos, y eso se logra con más muestras o con menor
        $f_s$. Subir $f_s$ con $N$ fijo amplía el ancho de banda analizado
        ($f_s/2$) pero acorta la ventana y empeora la resolución en la misma
        proporción (de \qty{15.6}{\hertz} a \qty{93.7}{\hertz} al pasar de
        \qty{8}{\kilo\sps} a \qty{48}{\kilo\sps} con $N = 512$). Para cada
        señal conviene la menor $f_s$ que cumpla con Nyquist y el mayor $N$
        que se pueda pagar.
  \item \textbf{La grilla de frecuencias se mueve con $f_s$ y con $N$.} Una
        misma senoidal puede caer justo en un coeficiente con una
        combinación (\qty{1}{\kilo\hertz} a \qty{8}{\kilo\sps}) y entre dos
        con otra (a \qty{22}{\kilo\sps} o \qty{48}{\kilo\sps}), con un error
        de amplitud de hasta \qty{3.92}{\deci\bel} y fuga hacia los
        coeficientes vecinos. Duplicar $N$ achica ese error de frecuencia a
        la mitad pero no lo elimina; para medir amplitudes con precisión hay
        que aplicar una ventana.
  \item \textbf{Una ventana más larga baja el piso de ruido por
        coeficiente.} El ruido de cuantización se reparte entre los $N$
        coeficientes mientras que la energía de una senoidal se concentra en
        uno, así que cada vez que se duplica $N$ el piso de ruido baja
        \qty{3}{\deci\bel} respecto del tono y se pueden ver componentes más
        débiles.
  \item \textbf{Una ventana más larga cuesta más en todo lo demás.} La
        memoria crece con $N$, el cálculo con $N\log_2 N$, el tiempo de
        transmisión con $N$ y la latencia (el tiempo hasta tener el primer
        coeficiente) con $T_v$. Además, una ventana larga promedia la señal
        durante más tiempo: si la señal cambia dentro de la ventana, el
        espectro mezcla los distintos estados. Es el compromiso entre
        resolución en frecuencia y resolución temporal.
  \item \textbf{La frecuencia de muestreo fija también el aliasing.} Sin
        filtro antialias, toda componente por encima de $f_s/2$ se pliega
        dentro de la banda y es indistinguible de una componente real;
        subir $f_s$ es la única defensa en este sistema.
  \item \textbf{En esta implementación el costo de la FFT no es el cuello de
        botella.} El cálculo dura del orden de un milisegundo, mientras que
        la trama tarda entre uno y cuatro segundos en salir por la UART. El
        período de refresco depende casi sólo de $N$ y casi nada de $f_s$
        (tabla~\ref{tab:tiempos}).
\end{itemize}

\completar{Comentar lo observado en las figuras~\ref{fig:gui-8k-512}
a~\ref{fig:gui-cuadrada} y en la tabla~\ref{tab:medido}: ancho del pico y
piso de ruido al aumentar $N$, desplazamiento del pico y fuga al aumentar
$f_s$, frecuencia aparente de la senoidal de \qty{5}{\kilo\hertz} y
amplitud de los armónicos de la cuadrada.}

\subsection{Decisiones de diseño}

\begin{itemize}[nosep]
  \item \textbf{Procesamiento por bloques con una bandera de entrega}: la
        rutina del ADC llena el buffer y el lazo principal lo procesa, sin
        duplicar buffers ni deshabilitar interrupciones.
  \item \textbf{Configuración copiada al comenzar cada bloque}: la FFT, la
        cantidad de filas y la $f_s$ informada corresponden siempre a las
        muestras, aunque el usuario cambie la selección en el medio.
  \item \textbf{Descarte del bloque parcial} al cambiar $f_s$ o al perder
        muestras, para que cada ventana sea continua y uniforme.
  \item \textbf{FFT en punto flotante} con las instancias precalculadas de
        CMSIS-DSP, aprovechando la FPU y evitando el escalado de la FFT en
        punto fijo.
  \item \textbf{Trama de texto con los bits IEEE-754 en hexadecimal}: exacta,
        sin soporte de \texttt{printf} de flotantes, con índice y cierre
        para detectar errores, y con la $f_s$ real en la cabecera.
  \item \textbf{Touchpad atendido durante la transmisión}, para que siga
        respondiendo aunque una trama tarde varios segundos.
  \item \textbf{VREF encendido explícitamente} antes de inicializar el ADC,
        ya que al quitar el DAC dejó de encenderse como efecto lateral.
\end{itemize}

\subsection{Mejoras posibles}\label{sec:mejoras}

\begin{itemize}[nosep]
  \item \textbf{Transmitir menos}: sólo los $N/2+1$ coeficientes no
        redundantes, en binario y a una velocidad mayor (el COM virtual del
        MCU-Link admite velocidades mayores que \num{115200}~bit/s),
        reduciría el tiempo de transmisión más de diez veces.
  \item \textbf{Usar \texttt{arm\_rfft\_fast\_f32()}}, que aprovecha que la
        entrada es real y calcula la mitad de operaciones con la mitad de
        memoria.
  \item \textbf{Aplicar una ventana} (Hann o Hamming, disponibles en
        CMSIS-DSP) antes de la FFT para bajar la fuga y el error de
        amplitud, a cambio de un lóbulo principal más ancho.
  \item \textbf{Adquisición continua con doble buffer y DMA}, para no
        perder muestras durante el procesamiento si el enlace lo permitiera,
        o para promediar varios espectros en la placa antes de enviar uno.
  \item \textbf{Antirrebote de las teclas}: un rebote de SW2 puede alternar
        el modo dos veces y dejarlo como estaba.
  \item \textbf{Barreras de memoria explícitas}: la entrega del bloque
        depende del orden entre la escritura de \texttt{buffer\_in} (y de
        la configuración del bloque) y la de \texttt{buffer\_listo}.
        Funciona porque el núcleo es único y \texttt{buffer\_listo} es
        \texttt{volatile}, pero declarar también \texttt{volatile} la
        configuración del bloque o agregar \texttt{\_\_DMB()} haría el
        código robusto frente a optimizaciones más agresivas del
        compilador.
\end{itemize}

%%%%%%%%%%%%%%%%%%%%%%%%%%%%%%%%%%%%%%%%%%%%%%%%%%%%%%%%%%%%%%%%%%%%%%%%%%%%%%%
\clearpage
\appendix
\section{Código fuente}\label{sec:codigo}
%%%%%%%%%%%%%%%%%%%%%%%%%%%%%%%%%%%%%%%%%%%%%%%%%%%%%%%%%%%%%%%%%%%%%%%%%%%%%%%

El proyecto completo de MCUXpresso (incluidos los archivos generados por
Config Tools, los drivers del SDK y CMSIS-DSP) está en la carpeta
\archivo{TP3/} del repositorio, y los scripts de MATLAB en
\archivo{TP3/matlab/}. El repositorio es:
\begin{center}
\url{https://github.com/Mati-Costamagna/procesamiento-digital-de-senales}
\end{center}

\end{document}
